\documentclass[11pt,a4paper]{article}
\usepackage[utf8]{inputenc}
\usepackage[T1]{fontenc}
\usepackage[czech]{babel}
\usepackage[margin=2.2cm]{geometry}
\usepackage{amsmath,amssymb,amsthm}
\usepackage{parskip}
\usepackage{enumitem}
\setlist{nosep,leftmargin=1.4em}
\usepackage{booktabs}
\usepackage{array}
\usepackage{xcolor}
\definecolor{linkblue}{RGB}{20,60,150}
\definecolor{defbg}{RGB}{240,244,252}
\definecolor{defframe}{RGB}{100,130,190}
\definecolor{markcol}{RGB}{176,84,0}
\definecolor{markbg}{RGB}{253,246,236}
\usepackage[most]{tcolorbox}
\newtcolorbox{defbox}[1][]{breakable,colback=defbg,colframe=defframe,boxrule=0.6pt,arc=1.5mm,left=2mm,right=2mm,top=1mm,bottom=1mm,title={#1},fonttitle=\bfseries}
\newtcolorbox{poznbox}[1][]{breakable,colback=markbg,colframe=markcol,boxrule=0.6pt,arc=1.5mm,left=2mm,right=2mm,top=1mm,bottom=1mm,title={#1},fonttitle=\bfseries}
\newtheorem{veta}{Věta}[section]
\theoremstyle{definition}
\newtheorem{definice}{Definice}[section]
\usepackage[colorlinks=true,linkcolor=linkblue,urlcolor=linkblue]{hyperref}
\usepackage{algorithm}
\usepackage{algpseudocode}
\usepackage{graphicx}
\graphicspath{{obrazky/mas/}}
% Obrázek ze slidů: nepoužívá float (dokument floaty nemá), jen centrovaný blok s~popiskem.
\newcommand{\slajd}[3]{%
  \begin{center}%
  \includegraphics[width=#2\textwidth]{#1}\\[1.2mm]%
  {\small\itshape #3}%
  \end{center}}

% Značka pro látku, která není ve slidech NAIL106 (mas-18en.pdf).
\newcommand{\mimo}{\textcolor{markcol}{\textsf{\footnotesize[$\blacklozenge$~nad rámec slidů]}}}
\newcommand{\mimoq}[1]{\textcolor{markcol}{\textsf{\footnotesize[$\blacklozenge$~nad rámec slidů: #1]}}}

\title{\textbf{Multiagentní systémy}\\[0.3em]\large Učební text ke státní závěrečné zkoušce}
\author{SZZ Umělá inteligence, MFF UK}
\date{srpen 2026}

\begin{document}
\maketitle

\vspace{-1.5em}

\begin{center}
\begin{minipage}{0.93\textwidth}
\small
Okruh \uv{Multiagentní systémy} magisterské SZZ Umělá inteligence. Kapitoly $=$ věty
oficiálního znění okruhu \textbf{jedna ku jedné} (mapovací tabulka níže). Zdroje: slidy
NAIL106 \emph{Multiagentní systémy} (Neruda, Pilát, MFF UK); M.~Wooldridge: \emph{An
Introduction to MultiAgent Systems} (2.~vyd., 2009); Y.~Shoham, K.~Leyton-Brown:
\emph{Multiagent Systems} (2009); S.~Russell, P.~Norvig: \emph{Artificial Intelligence:
A~Modern Approach}. Anglické termíny kurzívou v~závorce při prvním výskytu (zkoušející je
používá, literatura je anglická).
\end{minipage}
\end{center}

\begin{center}
\begin{minipage}{0.93\textwidth}
\begin{poznbox}[Rozsah textu a~značení]
\small
\begin{itemize}
\item Rozsah kapitol $\sim$ váha látky na přednášce: nejvíc architektury agentů, komunikace,
  kooperace a~teorie her.
\item \mimo{} $=$ látka \textbf{mimo slidy} přednášky, ponechaná, protože ji \emph{jmenuje
  okruh} nebo bez ní výklad nedává smysl; nabídka ke~škrtnutí. Neoznačené $=$ doloženo ve
  slidech.
\item Kapitoly \ref{sec:cesta} (hledání cesty) a~\ref{sec:uceni} (učení agentů): ve slidech
  nejsou (přednáška je uvádí mezi nepokrytými), okruh je vyjmenovává $\Rightarrow$ nutné
  umět; zpracovány stručně.
\end{itemize}
\end{poznbox}
\end{minipage}
\end{center}

\section*{Mapování okruhu na kapitoly}
\addcontentsline{toc}{section}{Mapování okruhu na kapitoly}

\begin{center}
\small
\begin{tabular}{p{0.55\textwidth} p{0.38\textwidth}}
\toprule
\textbf{Věta oficiálního znění okruhu} & \textbf{Kapitola v~tomto textu} \\
\midrule
Architektura autonomního agenta; mechanismus výběru akcí.
  & \ref{sec:architektura}~Architektura autonomního agenta a~mechanismus výběru akcí \\
Metody pro řízení agentů; symbolické a~konekcionistické reaktivní plánování, hybridní
přístupy.
  & \ref{sec:rizeni}~Metody pro řízení agentů \\
Problém hledání cesty.
  & \ref{sec:cesta}~Problém hledání cesty \\
Komunikace a~znalosti v~multiagentních systémech, ontologie, řečové akty, FIPA-ACL,
protokoly.
  & \ref{sec:komunikace}~Komunikace a~znalosti v~MAS \\
Distribuované řešení problémů, kooperace, Nashova ekvilibria, Paretova efektivita, alokace
zdrojů, aukce.
  & \ref{sec:kooperace}~Distribuované řešení problémů, kooperace, teorie her a~aukce \\
Metody pro učení agentů; zpětnovazební učení.
  & \ref{sec:uceni}~Učení agentů a~zpětnovazební učení \\
Metodologie návrhu, jazyky a~prostředí multiagentních systémů.
  & \ref{sec:metodologie}~Metodologie návrhu, jazyky a~prostředí MAS \\
\bottomrule
\end{tabular}
\end{center}

\tableofcontents
\newpage

%=====================================================================
\section{Architektura autonomního agenta a~mechanismus výběru akcí}
\label{sec:architektura}
%=====================================================================

\subsection{Agent a~multiagentní systém}

Multiagentní systémy (\emph{multiagent systems}, MAS) --- samostatná oblast informatiky od
90.~let. Hnací trendy:
\begin{itemize}
\item ubikvitnost výpočetní techniky; konektivita a~distribuovanost (výpočet jako interakce);
\item rostoucí složitost úloh; delegace úloh na software; tlak na snadnost použití;
\item paradigma \uv{vstup $\to$ výpočet $\to$ výstup $\to$ konec} nestačí pro software trvale
  zasazený do prostředí, které se mění nezávisle na něm.
\end{itemize}

\begin{defbox}[Agent --- pracovní definice]
\textbf{Russell, Norvig (AIMA):} cokoli, co \emph{vnímá} prostředí senzory a~\emph{působí} na
ně efektory.

\textbf{Pattie Maes:} výpočetní systém obývající složité dynamické prostředí, v~němž autonomně
vnímá a~jedná, a~tím plní cíle či úkoly, pro které byl navržen.

\textbf{Wooldridge, Jennings:} počítačový systém zasazený do prostředí, schopný v~něm
\emph{autonomního jednání} vedoucího ke splnění delegovaných cílů.
\end{defbox}

Jednotná definice neexistuje --- Franklin a~Graesser (1996, \emph{Is it an agent, or just
a~program?}): desítky nekompatibilních definic. Podstatné: agent $\ne$ procedura či objekt.

\begin{defbox}[Multiagentní systém]
Systém více agentů, kteří spolu \emph{interagují} --- nejčastěji výměnou zpráv po síti.
Agenti \emph{nemusí sdílet společný cíl}, mohou mít vlastní i~protichůdné zájmy $\Rightarrow$
nutnost studovat vyjednávání, koordinaci a~dynamické rozdělování zdrojů.
\end{defbox}

\textbf{Autonomie.} Člověk plně, metoda v~Javě vůbec, agent mezi: volí \emph{způsob} řešení
(podcíle), ne hlavní cíl --- ten je delegován.

\paragraph{Vymezení vůči příbuzným oborům} (pět pohledů na MAS):
\begin{itemize}
\item \textbf{Softwarové inženýrství} --- další stupeň abstrakce: strojový kód $\to$ objekty
  $\to$ distribuované objekty $\to$ agenti; bonmot \uv{objekty to dělají zadarmo, agenti proto,
  že chtějí}.
\item \textbf{Distribuované a~paralelní systémy} --- mutex, deadlock, konsensus vyřešeny; agenti
  jsou však \emph{autonomní}, koordinace se řeší až za běhu.
\item \textbf{Umělá inteligence} --- tradičně MAS $\subseteq$ UI; Russell a~Norvig naopak: cíl
  UI = návrh inteligentních agentů. MAS užívá techniky UI + \emph{sociální} aspekty, které
  klasická UI přehlíží.
\item \textbf{Teorie her a~ekonomie} --- formální aparát racionálního rozhodování více subjektů;
  matematická teorie her řeší \emph{existenci}, MAS \emph{výpočetní} aspekty.
\item \textbf{Sociologie} --- společnosti agentů: inspirace ano, nutnost ne. Opačná vazba
  pevnější: sociologie, ekonomie i~ekologie staví simulační modely na agentech
  (\emph{agent-based modelling}, ABM).
\end{itemize}

\subsection{Vlastnosti inteligentního agenta}

\begin{itemize}
\item \textbf{Reaktivita} (\emph{reactive}) --- vnímá prostředí, reaguje na změny v~rozumném
  (reálném) čase.
\item \textbf{Proaktivita} (\emph{proactive}) --- vlastní cíle, naplňuje je z~vlastní
  iniciativy, ne jen jako odpověď na podnět.
\item \textbf{Sociálnost} (\emph{social ability}) --- komunikuje s~jinými agenty i~s~lidmi.
\end{itemize}

Čistě reaktivní i~čistě cílové chování = snadné; obtížná je \textbf{rovnováha reaktivity
a~proaktivity}: ani slepě držet plán, který ztratil smysl, ani stále přehodnocovat cíle
a~nikam nedojít. Vrací se: odhodlání agenta (odst.~\ref{ssec:commitment}), hybridní
architektury (odst.~\ref{ssec:hybridni}).

\subsection{Intencionální systémy a~intencionální postoj}

Agentům přisuzujeme \emph{mentální stavy} (\uv{agent věří, že\dots}, \uv{agent chce\dots}).
\textbf{Intencionální systém} (Dennett) --- entita, jejíž chování lze predikovat přisouzením
přesvědčení (\emph{beliefs}), tužeb (\emph{desires}) a~racionality. Hierarchie: \emph{1.~řád}
= přesvědčení a~touhy; \emph{2.~řád} = přesvědčení a~touhy o~přesvědčeních a~touhách (svých
i~cizích); atd.

Tři Dennettovy \emph{postoje} (\emph{stances}) k~predikci chování:
\begin{itemize}
\item \emph{fyzikální} --- stačí fyzikální zákony (hozený kámen: hmotnost, rychlost,
  gravitace, ne touhy);
\item \emph{návrhový} (\emph{design stance}) --- neznám fyziku, znám účel (nastavím budík
  bez znalosti elektroniky);
\item \emph{intencionální} --- popis pomocí přesvědčení, tužeb a~racionality.
\end{itemize}

Shoham: \uv{Vypínač je velmi kooperativní agent, který vede proud, kdykoli věří, že si to
přejeme.} Konzistentní, ale absurdní --- \emph{existuje jednodušší, stejně přesný popis}.
Intencionální postoj se vyplatí, kde fyzikální či návrhový popis neexistuje nebo je
nepraktický. \textbf{Pro informatiku: abstrakce pro zvládnutí složitosti} --- pro mnoho
programátorů hlavní přínos MAS.

\subsection{Prostředí a~jeho vlastnosti}

Typ prostředí určuje chování agenta. Klasifikace (Russell, Norvig):

\begin{center}
\begin{tabular}{p{0.28\textwidth} p{0.64\textwidth}}
\toprule
\textbf{Dichotomie} & \textbf{Význam} \\
\midrule
plně / částečně pozorovatelné \newline (\emph{fully / partially observable})
  & senzory dávají úplný stav prostředí --- nebo jen část (skryté informace, šum, omezený
    dosah) \\
statické / dynamické \newline (\emph{static / dynamic})
  & prostředí se mění jen akcemi agenta --- nebo i~jinak (jiní agenti, přírodní děje);
    ve statickém si agent může \uv{v~klidu rozmyslet} tah \\
deterministické / nedeterministické
  & akce má právě jeden výsledek --- nebo více (se známým či neznámým rozdělením
    pravděpodobnosti = stochastické prostředí) \\
diskrétní / spojité
  & konečný počet stavů a~akcí --- nebo spojité veličiny (poloha, rychlost) \\
\bottomrule
\end{tabular}
\end{center}

\textbf{Nepříjemné konstatování:} zajímavá prostředí (reálný svět, internet) = spojitá,
nedeterministická, dynamická, částečně pozorovatelná $\Rightarrow$ pevná tabulka pravidel
nestačí.

\emph{Příklady.} \textbf{Termostat}: vjemy \uv{je zima} / \uv{teplota OK}, akce zapnout /
vypnout; přesto nedeterministické prostředí (někdo otevře dveře). \textbf{Softwarový démon}
(unixový \texttt{xbiff}): prostředí = operační systém, vjemy = systémová volání, akce =
softwarové operace.

\subsection{Abstraktní architektura agenta}

Předpoklad: diskrétní prostředí --- neomezuje, spojité lze diskretizovat s~libovolnou
přesností.

\begin{defbox}[Prostředí, agent, běh]
\textbf{Prostředí}: konečně mnoho diskrétních stavů $E = \{e, e', \dots\}$. \textbf{Agent}:
konečná množina akcí $A = \{a, a', \dots\}$.

\textbf{Běh} (\emph{run}) --- posloupnost střídajících se stavů prostředí a~akcí:
\[
r = e_0 \xrightarrow{a_0} e_1 \xrightarrow{a_1} e_2 \xrightarrow{a_2} \cdots
\xrightarrow{a_{n-1}} e_n .
\]
$R$ --- všechny konečné běhy (nad $E$ a~$A$); $R^A \subseteq R$ běhy končící akcí;
$R^E \subseteq R$ běhy končící stavem prostředí.
\end{defbox}

\begin{defbox}[Přechodová funkce prostředí (\emph{state transformer function})]
\[
\tau : R^A \to 2^E
\]
Běh končící akcí $\mapsto$ množina možných následujících stavů. Důsledky:
\begin{itemize}
\item \textbf{prostředí má historii} --- další stav závisí na celém běhu, ne jen na poslední akci;
\item \textbf{prostředí je nedeterministické} --- výsledkem je množina stavů, který nastane,
  dopředu nevíme.
\end{itemize}
$\tau(r) = \emptyset$ $\Rightarrow$ běh skončil. Předpoklad: všechny běhy skončí.

\textbf{Prostředí} $\mathit{Env} = \langle E, e_0, \tau\rangle$, $e_0$ počáteční stav.
\end{defbox}

\begin{defbox}[Agent a~systém]
\textbf{Agent} --- funkce z~běhů končících stavem prostředí do akcí:
\[
Ag : R^E \to A .
\]
Rozhoduje podle \emph{celé historie} a~\emph{deterministicky}. $\mathcal{AG}$ --- množina všech
agentů.

\textbf{Systém} --- dvojice $\langle Ag, \mathit{Env}\rangle$. Posloupnost $(e_0, a_0, e_1, a_1,
\dots)$ je \emph{běh agenta $Ag$ v~prostředí $\mathit{Env}$}, právě když
\begin{enumerate}
\item $e_0$ je počáteční stav $\mathit{Env}$,
\item $a_0 = Ag(e_0)$,
\item pro každé $u > 0$: $e_u \in \tau\big((e_0,a_0,\dots,a_{u-1})\big)$
  a~$a_u = Ag\big((e_0,a_0,\dots,e_u)\big)$.
\end{enumerate}
$R(Ag,\mathit{Env})$ --- množina všech běhů $Ag$ v~$\mathit{Env}$.
\end{defbox}

\begin{defbox}[Funkční ekvivalence agentů]
$Ag_1$, $Ag_2$ jsou \textbf{funkčně ekvivalentní vzhledem k~$\mathit{Env}$} $\iff$
$R(Ag_1,\mathit{Env}) = R(Ag_2,\mathit{Env})$.

\textbf{Prostě funkčně ekvivalentní} (\emph{simply functionally equivalent}) --- funkčně
ekvivalentní vzhledem ke \emph{všem} prostředím.
\end{defbox}

Nástroj pro srovnání \emph{vyjadřovací síly} architektur: implementace nepodstatná, vedou-li
ke stejným pozorovatelným běhům.

\subsection{Čistě reaktivní agent a~agent s~vnitřním stavem}

\begin{defbox}[Čistě reaktivní agent (\emph{purely reactive / tropistic agent})]
\[
Ag : E \to A
\]
Reaguje jen na aktuální stav prostředí, historii ignoruje.

\textbf{Vztah k~obecnému agentovi:} ke každému čistě reaktivnímu agentovi existuje funkčně
ekvivalentní standardní agent; opačně \emph{neplatí}.

\emph{Termostat:} $Ag(e) = \mathrm{off}$ pro $e = $ \uv{teplota OK}, jinak $Ag(e) = \mathrm{on}$.
\end{defbox}

\begin{defbox}[Agent s~vnitřním stavem]
$I$ --- vnitřní stavy agenta, $\mathit{Per}$ --- vjemy. Agent = trojice funkcí:
\begin{align*}
\mathit{see} &: E \to \mathit{Per} && \text{(vnímání)}\\
\mathit{next} &: I \times \mathit{Per} \to I && \text{(aktualizace vnitřního stavu)}\\
\mathit{action} &: I \to A && \text{(výběr akce)}
\end{align*}
\slajd{mas-agent-stav.png}{0.42}{Agent s~vnitřním stavem za práce: start ve stavu $i_0$,
\emph{see} převede stav prostředí na vjem, \emph{next} z~vjemu a~starého stavu spočítá nový
stav, \emph{action} z~něj vybere akci --- ta zpětně mění prostředí.}

\textbf{Věta (bez důkazu):} ke každému agentovi s~vnitřním stavem existuje funkčně ekvivalentní
standardní agent. Vnitřní stav \emph{nezvyšuje vyjadřovací sílu}, je zásadní pro
\emph{efektivitu} a~přehlednost implementace.
\end{defbox}

$\mathit{see}$ formalizuje \textbf{pozorovatelnost}: $e_1 \ne e_2$ se stejným vjemem = agent je
nerozliší (částečná pozorovatelnost). Krajní případy: $|\mathit{Per}| = |E|$, $\mathit{see}$
prostá (plně pozorovatelné) vs.\ $|\mathit{Per}| = 1$ (agent nevidí nic).

\subsection{Utilita: jak agentovi říct, co má dělat}

Ne zadrátované chování: agentovi říkáme \emph{co}, ne \emph{jak}. Standardní postup UI: cíl
\emph{nepřímo}, mírou úspěšnosti.

\begin{defbox}[Utilita]
\textbf{Utilita stavu:} $u : E \to \mathbb{R}$; cíl agenta = stavy s~vysokou utilitou.

Utilita stavů je krátkozraká $\Rightarrow$ \textbf{utilita běhu}
\[
u : R \to \mathbb{R} ,
\]
pro dlouhodobě běžící agenty. Ze stavů např.\ utilita \emph{nejhoršího} navštíveného stavu
nebo \emph{průměr} navštívených stavů --- podle úlohy.
\end{defbox}

\begin{defbox}[Optimální agent]
$P(r \mid Ag, \mathit{Env})$ --- pravděpodobnost běhu $r$ agenta $Ag$ v~$\mathit{Env}$,
$\sum_{r \in R(Ag,\mathit{Env})} P(r \mid Ag,\mathit{Env}) = 1$. Optimální agent maximalizuje
\emph{očekávanou utilitu}:
\[
Ag_{\mathrm{opt}} = \arg\max_{Ag \in \mathcal{AG}} \sum_{r \in R(Ag,\mathit{Env})}
u(r)\, P(r \mid Ag, \mathit{Env}).
\]
\end{defbox}

Elegantní, ale \textbf{nekonstruktivní} --- neříká, jak agenta navrhnout; utilitu bývá těžké
zadat. (Peníze nejsou dobrá utilita člověka: nelineární, jiná pro bohaté a~chudé.)

\textbf{Příklad: Tileworld} --- klasické testovací prostředí architektur. Šachovnice, pohyb do
čtyř směrů; díry, překážky, bloky; cíl: zatlačit bloky do děr. \emph{Dynamické}: díry, překážky
i~bloky náhodně vznikají a~zanikají. Utilita běhu
\[
u(r) = \frac{N_f}{N_a},
\]
$N_f$ = díry zaplněné agentem během $r$, $N_a$ = všechny díry vzniklé během $r$. Nutné
\emph{reagovat} (tlačený blok zmizel) i~\emph{využívat příležitosti} (vedle vznikl nový blok).
Historicky nástroj pro kompromis reaktivita vs.\ proaktivita (odst.~\ref{ssec:commitment}).

\subsection{Predikátová specifikace úlohy}

Utilita do $\{0,1\}$: běh úspěšný ($u(r)=1$), nebo ne.

\begin{defbox}[Prostředí úlohy (\emph{task environment})]
\textbf{Predikátová specifikace} $\Psi : R \to \{0,1\}$: $\Psi(r)$ platí $\iff$ $u(r) = 1$.
\textbf{Prostředí úlohy} --- dvojice $\langle \mathit{Env}, \Psi \rangle$: vlastnosti systému
(přes $\mathit{Env}$) + kritérium splnění úlohy (přes $\Psi$).

$R_\Psi(Ag,\mathit{Env}) = \{ r \in R(Ag,\mathit{Env}) \mid \Psi(r) \}$.
\end{defbox}

Kdy agent úlohu \emph{řeší} --- tři výklady:
\begin{itemize}
\item \textbf{Pesimistický:} $R_\Psi(Ag,\mathit{Env}) = R(Ag,\mathit{Env})$ --- \emph{všechny}
  běhy splňují $\Psi$.
\item \textbf{Optimistický:} $\exists r \in R(Ag,\mathit{Env}): \Psi(r)$ --- \emph{alespoň jeden}.
\item \textbf{Realistický:} $\tau$ rozšířena o~rozdělení pravděpodobnosti, úspěch měřen
  \[ P(\Psi \mid Ag,\mathit{Env}) = \sum_{r \in R_\Psi(Ag,\mathit{Env})} P(r \mid Ag,\mathit{Env}). \]
\end{itemize}

\begin{defbox}[Typy úloh]
\textbf{Úlohy dosažení} (\emph{achievement tasks}) --- \uv{něco najít}: cílové stavy
$G \subseteq E$; $\Psi(r)$ platí, obsahuje-li $r$ stav z~$G$. Typicky klasické úlohy UI
(prohledávání).

\textbf{Úlohy udržení} (\emph{maintenance tasks}) --- \uv{něčemu se vyhnout}: \uv{špatné} stavy
$B \subseteq E$; $\Psi(r)$ neplatí, obsahuje-li $r$ stav z~$B$. Typicky hry: $B$ = prohrané
pozice, prostředí = soupeř.

\textbf{Kombinace:} dosáhnout $G$ a~vyhnout se $B$.
\end{defbox}

\subsection{Mechanismus výběru akcí --- přehled}
\label{ssec:vyber-akci}

\begin{defbox}[Mechanismus výběru akcí (\emph{action selection mechanism})]
Procedura volby akce v~každém okamžiku podle vjemů a~vnitřního stavu.
\textbf{Architektura agenta} --- softwarová architektura realizující mechanismus; Maes (1991): \uv{konkrétní metodologie stavby agentů, která určuje, jak lze agenta rozložit na
množinu komponent a~jak mají tyto komponenty interagovat, aby dohromady odpověděly na otázku,
jak data ze senzorů a~aktuální vnitřní stav určují akce a~budoucí vnitřní stav agenta.}
\end{defbox}

Výběr akce = \textbf{klíčový problém návrhu agenta}; kapitola~\ref{sec:rizeni} = přehled
realizací:

\begin{center}
\small
\begin{tabular}{p{0.20\textwidth} p{0.34\textwidth} p{0.38\textwidth}}
\toprule
\textbf{Rodina} & \textbf{Mechanismus výběru akce} & \textbf{Typický zástupce} \\
\midrule
symbolická / deliberativní
  & logická dedukce nad databází formulí; plánování
  & deduktivní agent, STRIPS, BDI, PRS \\
reaktivní (symbolická)
  & pravidla \uv{situace $\to$ akce} s~pevnou prioritou
  & subsumpční architektura (Brooks) \\
reaktivní (konekcionistická)
  & šíření aktivace v~síti modulů, vyhrává nejaktivnější
  & agent network architecture (Maes) \\
hybridní
  & vrstvy různé abstrakce + arbitráž mezi nimi
  & TouringMachines, InteRRaP \\
koaliční
  & soutěž dynamicky vznikajících koalic procesů o~\uv{vědomí}
  & IDA / globální pracovní prostor \\
naučená
  & maximalizace naučené $Q$-funkce / stochastická politika
  & Q-learning, actor-critic (kap.~\ref{sec:uceni}) \\
\bottomrule
\end{tabular}
\end{center}

\subsection{Shrnutí ke~zkoušce}

\begin{itemize}
\item Agent = autonomní systém zasazený do prostředí, vnímá a~jedná, plní delegované cíle.
  MAS = interagující agenti \emph{bez nutně sdíleného cíle}.
\item Tři vlastnosti: reaktivita, proaktivita, sociálnost; těžké je vyvážení.
\item Dennett: intencionální systém, tři postoje (fyzikální, návrhový, intencionální);
  intencionální postoj = abstrakce pro zvládnutí složitosti.
\item Prostředí: 4 dichotomie (pozorovatelnost, statické/dynamické, determinismus,
  diskrétnost); reálná prostředí v~nejtěžší variantě.
\item Formalismus: $E$, $A$, běh $r = e_0 a_0 e_1 \dots e_n$, $R^A$/$R^E$,
  $\tau : R^A \to 2^E$ (historie + nedeterminismus),
  $\mathit{Env}=\langle E,e_0,\tau\rangle$, $Ag : R^E \to A$, systém $\langle Ag,\mathit{Env}\rangle$.
\item Funkční ekvivalence (vzhledem k~prostředí / prostá). Čistě reaktivní $Ag: E \to A$
  (slabší); agent s~vnitřním stavem ($\mathit{see}$, $\mathit{next}$, $\mathit{action}$) ---
  oba převeditelné na standardního agenta.
\item Utilita $u : E \to \mathbb{R}$ vs.\ $u : R \to \mathbb{R}$; optimální agent maximalizuje
  \emph{očekávanou} utilitu --- nekonstruktivní. Tileworld, $u(r)=N_f/N_a$.
\item Prostředí úlohy $\langle \mathit{Env}, \Psi\rangle$; pesimista/optimista/realista;
  úlohy dosažení ($G$) a~udržení ($B$).
\item \textbf{Mechanismus výběru akcí} = jádro otázky: vyjmenovat rodiny mechanismů
  a~zástupce (tabulka výše).
\end{itemize}

%=====================================================================
\section{Metody pro řízení agentů: symbolické a~konekcionistické reaktivní
plánování, hybridní přístupy}
\label{sec:rizeni}
%=====================================================================

Jádro okruhu: konkrétní realizace mechanismu výběru akcí z~odst.~\ref{ssec:vyber-akci},
od symbolického (deliberativního) pólu přes reaktivní k~hybridním architekturám:

\begin{center}
\small
\begin{tabular}{@{}p{0.26\textwidth} p{0.66\textwidth}@{}}
\toprule
\textbf{Přístup} & \textbf{Kde je v~této kapitole} \\
\midrule
symbolické (deliberativní) řízení
  & deduktivní agent, BDI, STRIPS, PRS (odst.\ 2.1--2.6) \\
symbolické reaktivní plánování
  & subsumpční architektura (Brooks), odst.~\ref{ssec:subsumpce} \\
konekcionistické reaktivní plánování
  & agent network architecture (Maes), odst.~\ref{ssec:maes} \\
hybridní přístupy
  & vrstvené architektury, TouringMachines, InteRRaP, IDA, odst.~\ref{ssec:hybridni} a~dále \\
\bottomrule
\end{tabular}
\end{center}
\subsection{Klasický přístup UI a~deduktivní agent}

Dva pilíře klasické UI:
\begin{enumerate}
\item \textbf{Symbolická reprezentace} prostředí a~chování --- logické formule.
\item \textbf{Syntaktická manipulace} s~ní --- dedukce, dokazování vět.
\end{enumerate}
Teorie chování agenta $=$ program (\emph{spustitelná specifikace}). Dva klasické problémy:
\begin{itemize}
\item \textbf{Problém transdukce} (\emph{transduction problem}) --- převod reálného světa na
  symboly (počítačové vidění, zpracování přirozeného jazyka, učení).
\item \textbf{Problém reprezentace a~uvažování} --- jak znalosti reprezentovat a~efektivně nad
  nimi odvozovat (dokazovače, plánovače).
\end{itemize}

\begin{defbox}[Deduktivní agent (agent jako dokazovač vět)]
$L$ --- formule predikátové logiky prvního řádu, $D = 2^L$ --- databáze formulí; vnitřní stav
agenta $\mathit{DB} \in D$. Uvažování $=$ odvozovací pravidla $\rho$:
\[
\mathit{DB} \vdash_\rho \varphi \quad\text{---\ formuli } \varphi \text{ lze odvodit z } \mathit{DB}
\text{ pomocí pravidel } \rho .
\]
Agent $=$ trojice $\mathit{see} : E \to \mathit{Per}$, $\mathit{next} : D \times \mathit{Per} \to D$,
$\mathit{action} : D \to A$; $\mathit{action}$ dána procedurou níže.
\end{defbox}

\noindent\textbf{Výběr akce jako dokazování.}
\begin{algorithmic}[1]
\Function{action}{$\mathit{DB} : D$} : $A$
  \ForAll{$a \in A$}
    \If{$\mathit{DB} \vdash_\rho \mathit{Do}(a)$} \Return $a$ \EndIf
  \EndFor
  \ForAll{$a \in A$}
    \If{$\mathit{DB} \not\vdash_\rho \neg \mathit{Do}(a)$} \Return $a$ \EndIf
  \EndFor
  \State \Return $\mathit{null}$
\EndFunction
\end{algorithmic}

1.~cyklus: akce \emph{dokazatelně} žádoucí ($\mathit{Do}(a)$); 2.~cyklus: akce aspoň
\emph{konzistentní} s~DB (nelze dokázat $\neg\mathit{Do}(a)$); jinak nic.

\medskip
\noindent\textbf{Příklad: vysavačový svět 3$\times$3.}
Predikáty $\mathit{In}(x,y)$ (pozice agenta), $\mathit{Dirt}(x,y)$ (špína), $\mathit{Facing}(d)$
(směr). Pravidla:
\begin{align*}
\mathit{In}(x,y) \wedge \mathit{Dirt}(x,y) &\Rightarrow \mathit{Do}(\mathit{suck}) \\
\mathit{In}(0,0) \wedge \mathit{Facing}(\mathit{north}) \wedge \neg\mathit{Dirt}(0,0)
  &\Rightarrow \mathit{Do}(\mathit{forward}) \\
\mathit{In}(0,2) \wedge \mathit{Facing}(\mathit{north}) \wedge \neg\mathit{Dirt}(0,2)
  &\Rightarrow \mathit{Do}(\mathit{turn}) \qquad \dots
\end{align*}
Agent projde svět \uv{hadem} 00--01--02--12--11--10--\dots; úklid má přednost před pohybem
(pravidlo $\mathit{suck}$ první, ostatní mají v~předpokladu $\neg\mathit{Dirt}$).

\medskip
\noindent\textbf{Klady a~zápory.}
\begin{itemize}
\item[$+$] Elegantní, jasná sémantika; specifikace $=$ program.
\item[$-$] Prakticky nepoužitelné pro netriviální úlohy: dokazování pomalé, nemusí skončit.
\item[$-$] \textbf{Problém vypočitatelné racionality} (\emph{calculative rationality}): akce
  optimální pro stav prostředí \emph{na začátku} odvozování --- mezitím se změnilo. Katastrofa
  v~rychle se měnícím prostředí.
\item[$-$] Těžký návrh $\mathit{see}$ (obrázek $\to$ formule, reprezentace času).
\end{itemize}
Prvky přístupu přežívají v~pozdějších architekturách, zejména BDI.

\subsection{Praktické uvažování: architektura BDI}

\begin{defbox}[Praktické uvažování (\emph{practical reasoning})]
Inspirace lidským rozhodováním. \emph{Teoretické} uvažování (Sókratés) $\to$ co si o~světě
myslíme; \emph{praktické} $\to$ \textbf{akce}. Dvě fáze:
\begin{itemize}
\item \textbf{Deliberace} (\emph{deliberation}) --- \emph{čeho} chceme dosáhnout
  (\uv{dodělat magisterské studium}).
\item \textbf{Uvažování o~prostředcích} (\emph{means-ends reasoning}) $=$ plánování ---
  \emph{jak} (plán, jak dostudovat).
\end{itemize}
Obojí nesmí trvat příliš dlouho.
\end{defbox}

\begin{defbox}[Beliefs --- Desires --- Intentions]
\textbf{Beliefs} (přesvědčení) --- informační stav, znalosti o~světě. Ne \uv{znalosti}:
jsou \emph{subjektivní}, \emph{nemusí být pravdivé}, \emph{mohou se změnit}. Mohou obsahovat
i~odvozovací pravidla (dopředné řetězení).

\textbf{Desires} (touhy) --- motivační stav, cíle/situace k~dosažení. \emph{Mohou být vzájemně
nekonzistentní} (učit se i~jít na pivo).

\textbf{Intentions} (záměry) --- stavy světa, k~jejichž dosažení se agent \emph{zavázal}; vedou
k~akcím (agent kvůli nim plánuje). \textbf{Přetrvávají}, dokud
\begin{enumerate}
\item jich agent nedosáhne, nebo
\item nezačne věřit, že jsou nedosažitelné, nebo
\item nepominou důvody, které k~nim vedly.
\end{enumerate}
Dále \emph{omezují další deliberaci} (nezvažuji možnosti neslučitelné se záměrem)
a~\emph{ovlivňují budoucí přesvědčení} (počítám s~úspěchem záměru).
\end{defbox}

\noindent\textbf{Touha vs.\ záměr} (Bratman, 1990): touha (zahrát si odpoledne basketbal) $=$
potenciální vliv soupeřící s~ostatními touhami; záměr $=$ rozhodnuto, pro a~proti už nevažuji.

\begin{defbox}[Funkce deliberace]
$B \subseteq \mathit{Bel}$, $D \subseteq \mathit{Des}$, $I \subseteq \mathit{Int}$ --- aktuální
přesvědčení, touhy, záměry. Tři funkce:
\begin{itemize}
\item $\mathit{brf} : 2^{\mathit{Bel}} \times \mathit{Per} \to 2^{\mathit{Bel}}$ ---
  \emph{belief revision function}, aktualizace přesvědčení podle vjemu;
\item $\mathit{options} : 2^{\mathit{Bel}} \times 2^{\mathit{Int}} \to 2^{\mathit{Des}}$ ---
  generování možností (tužeb) z~přesvědčení a~stávajících záměrů;
\item $\mathit{filter} : 2^{\mathit{Bel}} \times 2^{\mathit{Des}} \times 2^{\mathit{Int}}
  \to 2^{\mathit{Int}}$ --- výběr záměrů z~možností.
\end{itemize}
Klíčové rozdělení: \emph{generování} levné a~široké, \emph{filtrování} $=$ vlastní rozhodování
(bere v~úvahu stávající závazky).
\end{defbox}

\subsection{Plánování: STRIPS}

\begin{defbox}[Uvažování o~prostředcích / plánování]
Jak dosáhnout cíle (záměru) dostupnými prostředky (akcemi).
\begin{itemize}
\item \textbf{Vstup:} cíl ($=$ záměr $I$); stav prostředí ($=$ přesvědčení $B$); akce
  $\mathit{Ac}$.
\item \textbf{Výstup:} plán $=$ posloupnost akcí, po níž je cíl splněn.
\end{itemize}
\end{defbox}

\textbf{STRIPS} (Fikes, Nilsson, 1971): svět $=$ množina formulí FOL; akce $=$ \emph{schéma}
s~předpoklady a~efekty (přidávané/mazané fakty). Algoritmus: rozdíl cíl vs.\ stav, vhodná akce
ho zmenší (\emph{means-ends analysis}). Hezké, málo praktické --- utápí se v~detailech nízké úrovně.

\begin{defbox}[Deskriptor akce, plánovací problém, plán]
\textbf{Deskriptor akce} $a$: trojice $\langle P_a, D_a, A_a\rangle$ ---
předpoklady (\emph{preconditions}), mazaná množina (\emph{delete list}), přidávaná množina
(\emph{add list}); uzemněné atomické formule (bez spojek a~proměnných).

\textbf{Plánovací problém}: $\langle B_0, O, G\rangle$ --- počáteční přesvědčení, deskriptory
všech akcí $O = \{\langle P_a,D_a,A_a\rangle : a \in \mathit{Ac}\}$, cíl $G$ (množina formulí).

\textbf{Plán} $\pi = (a_1,\dots,a_n)$, $a_i \in \mathit{Ac}$, určuje databáze
\[
B_i = (B_{i-1} \setminus D_{a_i}) \cup A_{a_i}, \qquad i = 1,\dots,n .
\]
\textbf{Přípustný} (\emph{admissible}): $B_{i-1} \models P_{a_i}$ pro každé $i$.
\textbf{Korektní/validní} (\emph{correct}): přípustný a~$B_n \models G$.

\textbf{Plánování} $=$ pro $\langle B_0,O,G\rangle$ najít korektní plán, nebo prohlásit, že neexistuje.
\end{defbox}

\medskip
\noindent\textbf{Příklad: svět kostek (\emph{blocks world}).}
Predikáty: $\mathit{On}(x,y)$ ($x$ na $y$), $\mathit{OnTable}(x)$, $\mathit{Clear}(x)$
(na $x$ nic), $\mathit{Holding}(x)$ (rameno drží $x$), $\mathit{ArmEmpty}$.
Akce:
\begin{center}
\small
\begin{tabular}{p{0.20\textwidth} p{0.36\textwidth} p{0.36\textwidth}}
\toprule
\textbf{Akce} & \textbf{Pre / Del} & \textbf{Add} \\
\midrule
$\mathit{Stack}(x,y)$
  & Pre: $\{\mathit{Clear}(y), \mathit{Holding}(x)\}$ \newline
    Del: $\{\mathit{Clear}(y), \mathit{Holding}(x)\}$
  & $\{\mathit{ArmEmpty}, \mathit{On}(x,y)\}$ \\
$\mathit{UnStack}(x,y)$
  & Pre: $\{\mathit{On}(x,y),$ $\mathit{Clear}(x),\mathit{ArmEmpty}\}$ \newline
    Del: $\{\mathit{On}(x,y),\mathit{ArmEmpty}\}$
  & $\{\mathit{Holding}(x), \mathit{Clear}(y)\}$ \\
$\mathit{Pickup}(x)$
  & Pre: $\{\mathit{OnTable}(x),$ $\mathit{Clear}(x),\mathit{ArmEmpty}\}$ \newline
    Del: $\{\mathit{OnTable}(x),\mathit{ArmEmpty}\}$
  & $\{\mathit{Holding}(x)\}$ \\
$\mathit{PutDown}(x)$
  & Pre: $\{\mathit{Holding}(x)\}$ \newline
    Del: $\{\mathit{Holding}(x)\}$
  & $\{\mathit{ArmEmpty}, \mathit{OnTable}(x)\}$ \\
\bottomrule
\end{tabular}
\end{center}
\begin{align*}
B_0 &= \{\mathit{Clear}(A), \mathit{On}(A,B), \mathit{OnTable}(B),
\mathit{OnTable}(C), \mathit{Clear}(C), \mathit{ArmEmpty}\},\\
G &= \{\mathit{OnTable}(A), \mathit{OnTable}(B), \mathit{OnTable}(C)\}.
\end{align*}
Korektní plán $\pi = (\mathit{UnStack}(A,B),\ \mathit{PutDown}(A))$: $B_0 \models$ Pre
1.~akce;
\[ B_1 = (B_0 \setminus \{\mathit{On}(A,B),\mathit{ArmEmpty}\}) \cup
\{\mathit{Holding}(A),\mathit{Clear}(B)\} \models \mathit{Holding}(A); \]
$B_2 \models G$ $\Rightarrow$ přípustný i~korektní.

\begin{defbox}[Pomocné funkce nad plány]
$\mathit{pre}(\pi)$ (předpoklad), $\mathit{body}(\pi)$ (tělo), $\mathit{empty}(\pi)$,
$\mathit{execute}(\pi)$, $\mathit{hd}(\pi)$ (první akce), $\mathit{tail}(\pi)$ (zbytek),
$\mathit{sound}(\pi,I,B)$ (korektnost vzhledem k~záměrům $I$ a~přesvědčením $B$);
plánovací funkce
$\mathit{plan} : 2^{\mathit{Bel}} \times 2^{\mathit{Int}} \times 2^{\mathit{Ac}} \to \mathit{Plan}$.
\end{defbox}

\textbf{Knihovna plánů.} Online konstrukce plánů je drahá; $\mathit{plan}()$ často $=$
\emph{knihovna hotových plánů}: projít, zkontrolovat předpoklady vs.\ přesvědčení a~efekty vs.\
cíl. Jádro PRS (sekce~\ref{ssec:prs}).

\subsection{Implementace BDI agenta}

\noindent\textbf{Hlavní řídicí smyčka BDI agenta:}
\begin{algorithmic}[1]
\State $B \gets B_0;\quad I \gets I_0$
\While{$\mathbf{true}$}
  \State $v \gets \mathit{see}()$
  \State $B \gets \mathit{brf}(B,v)$
  \State $D \gets \mathit{options}(B,I)$
  \State $I \gets \mathit{filter}(B,D,I)$
  \State $\pi \gets \mathit{plan}(B,I,\mathit{Ac})$
  \While{$\neg\big(\mathit{empty}(\pi) \vee \mathit{succeeded}(I,B) \vee \mathit{impossible}(I,B)\big)$}
    \State $a \gets \mathit{hd}(\pi);\quad \mathit{execute}(a);\quad \pi \gets \mathit{tail}(\pi)$
    \State $v \gets \mathit{see}();\quad B \gets \mathit{brf}(B,v)$
    \If{$\mathit{reconsider}(I,B)$}
      \State $D \gets \mathit{options}(B,I);\quad I \gets \mathit{filter}(B,D,I)$
    \EndIf
    \If{$\neg\,\mathit{sound}(\pi,I,B)$}
      \State $\pi \gets \mathit{plan}(B,I,\mathit{Ac})$ \Comment{přeplánování}
    \EndIf
  \EndWhile
\EndWhile
\end{algorithmic}

$\mathit{succeeded}(I,B)$ --- záměry $I$ platí za přesvědčení $B$;
$\mathit{impossible}(I,B)$ --- nemohou platit. Vnitřní cyklus končí: plán vyčerpán, cíl splněn,
nebo nedosažitelný.

\subsection{Odhodlání agenta (\emph{commitment})}
\label{ssec:commitment}

\begin{defbox}[Strategie odhodlání vůči záměrům]
Kdy opustit cíl?
\begin{itemize}
\item \textbf{Slepé odhodlání} (\emph{blind commitment}, též \emph{fanatical}) --- dokud
  \emph{nezačne věřit, že ho dosáhl}. (Nikdy se nevzdá.)
\item \textbf{Jednosměrné odhodlání} (\emph{single-minded commitment}) --- dokud nezačne věřit,
  že ho dosáhl, \emph{nebo} že už není dosažitelný.
\item \textbf{Otevřené odhodlání} (\emph{open-minded commitment}) --- dokud věří, že je
  dosažitelný (a~stále chtěný) --- lze opustit i~při změně motivace.
\end{itemize}
\end{defbox}

\textbf{Odhodlání vůči plánu} vs.\ \textbf{vůči cíli}: algoritmus výše je odhodlán plánu ---
opustí ho při splnění cíle, nedosažitelnosti, nebo prázdném plánu. Těžší: kdy
\emph{přehodnotit sám záměr}? Dilema --- deliberace \textbf{trvá čas}, prostředí se mění:
\begin{itemize}
\item \emph{nepřehodnocuje} $\to$ upne se k~nedosažitelným cílům;
\item přehodnocuje \emph{příliš často} $\to$ samé uvažování, nic nestihne.
\end{itemize}

\textbf{Metařízení} (Wooldridge, Parsons, 1999): $\mathit{reconsider}()$ \emph{výpočetně
jednodušší} než samo přehodnocení. Kvalita: dobrá $\iff$ kdykoli navrhne přehodnocení, agent
po deliberaci záměr \emph{skutečně} změní (a~naopak).

\begin{defbox}[Odvážný vs.\ opatrný agent]
\textbf{Odvážný agent} (\emph{bold agent}) --- přehodnocuje až po dokončení celého plánu (nikdy
plán nepřeruší kvůli deliberaci).

\textbf{Opatrný agent} (\emph{cautious agent}) --- přehodnocuje po každé akci plánu.

\textbf{Míra odvahy} (\emph{level of boldness}) --- počet akcí mezi dvěma deliberacemi.

\textbf{Dynamičnost prostředí} (\emph{rate of world change}) --- počet změn prostředí za jeden
cyklus agenta.

\textbf{Efektivita agenta} $=$ dosažené záměry $/$ všechny záměry.
\end{defbox}

\textbf{Experiment (Kinny, Georgeff, Tileworld):}
\begin{itemize}
\item svět se mění \emph{pomalu} $\to$ efektivnější \textbf{odvážní} (deliberace $=$ zbytečná režie);
\item \emph{rychle} $\to$ \textbf{opatrní} (plány rychle zastarávají).
\end{itemize}
Kvantitativní podoba kompromisu reaktivita vs.\ proaktivita.

\subsection{Procedurální odvozovací systém (PRS)}
\label{ssec:prs}

\begin{defbox}[PRS --- \emph{Procedural Reasoning System}]
Georgeff a~kol., 80.~léta, Stanford. \textbf{První implementace BDI}, jedna z~prakticky
nejúspěšnějších architektur; reimplementace: \textbf{AgentSpeak/Jason}, \textbf{JAM},
\textbf{JACK}, \textbf{Jadex}. Nasazení: OASIS (řízení letového provozu v~Sydney), SPOC, SWARMM.
\end{defbox}

\slajd{mas-prs.png}{0.46}{Architektura PRS agenta. Data ze senzorů vstupují do databáze
\emph{beliefs} (atomy FOL prologovského typu); \emph{interpret} propojuje beliefs, \emph{desires}
(cíle), \emph{knihovnu plánů} a~zásobník \emph{intentions} a~vydává akce.}

\begin{defbox}[Plán v~PRS]
Knihovna hotových plánů $=$ \emph{procedurální znalosti} agenta. \textbf{Neplánuje se od nuly},
jen se \emph{vybírá}. Tři části plánu:
\begin{itemize}
\item \textbf{Goal} (\emph{invocation condition} / \emph{cue}) --- co má platit po vykonání;
  na jaký cíl se plán hodí.
\item \textbf{Context} --- podmínka spuštění (konzistentní s~přesvědčeními).
\item \textbf{Body} --- akce k~vykonání.
\end{itemize}
Tělo \emph{není} jen lineární posloupnost akcí --- může obsahovat \emph{cíle}:
\uv{dosáhni $f$}, \uv{dosáhni $f$ nebo $g$}, \uv{dosahuj $f$, dokud nenastane $g$}
$\Rightarrow$ hierarchický rozklad cílů na podcíle.
\end{defbox}

\textbf{Běh interpretu:}
\begin{enumerate}
\item Inicializace: $B_0$ a~cíl nejvyšší úrovně (\emph{top-level goal}).
\item Zásobník záměrů: aktuální cíle v~různém stavu rozpracovanosti.
\item Záměr z~vrcholu zásobníku $\to$ v~knihovně hledá plány s~odpovídajícím cílem.
\item Z~nich ty s~\emph{kontextem} konzistentním s~přesvědčeními $=$ aktuální
  \emph{options} / \emph{desires}.
\item \textbf{Deliberace} $=$ výběr jednoho záměru:
  \begin{itemize}
  \item původní PRS: \emph{metaplány} (plány o~plánech modifikující záměry) --- příliš složité;
  \item praktičtější: číselné \emph{ohodnocení} plánu (očekávaná utilita), vybrat nejvyšší.
  \end{itemize}
\item Vykonání plánu $\to$ může přidat další záměry na zásobník.
\item Selhání plánu $\to$ jiný záměr z~možností.
\end{enumerate}

\medskip
\noindent\textbf{Ukázka --- Jason/AgentSpeak} (agent nosí majiteli pivo z~ledničky):
\begin{quote}\small\ttfamily
\noindent available(beer,fridge).\ \ limit(beer,10).\ \ \ \ /* počáteční beliefs */\\
too\_much(B) :- .count(consumed(\_,\_,\_,B),Q) \& limit(B,L) \& Q > L.\ \ /* pravidlo */\\[0.3em]
+!has(owner,beer)\ \ \ \ \ \ \ \ \ \ \ \ \ \ \ \ \ \ \ \ \ \ \ /* spouštěcí událost */\\
\hspace*{1em}: available(beer,fridge) \& not too\_much(beer)\ \ \ \ /* kontext */\\
\hspace*{1em}<- !at(robot,fridge); open(fridge); get(beer); close(fridge);\ \ /* tělo */\\
\hspace*{2em} !at(robot,owner); hand\_in(beer).\\[0.3em]
+!at(robot,P) : not at(robot,P) <- move\_towards(P); !at(robot,P).
\end{quote}
\texttt{+!g : c <- b} $=$ \uv{při události přidání cíle \texttt{g} a~platném kontextu \texttt{c}
vykonej tělo \texttt{b}}. \texttt{!} --- \emph{achievement goal} (dosáhni);
\texttt{+}/\texttt{-} --- přidání/odebrání přesvědčení nebo cíle; \texttt{.send}, \texttt{.count}
--- vestavěné akce.

\subsection{Reaktivní přístup a~jeho východiska}

80.\ a~90.\ léta: drobné úpravy symbolického přístupu nestačí $\to$ alternativní paradigmata UI.
Společný jmenovatel: \textbf{odmítnutí symbolické reprezentace} a~uvažování syntaktickou
manipulací. Tři teze:

\begin{defbox}[Tři principy reaktivního přístupu]
\textbf{Interakce} (\emph{situatedness}) --- inteligentní chování závisí na prostředí, do něhož
je agent \emph{zasazen}; nelze studovat izolovaně.

\textbf{Ztělesnění} (\emph{embodiment}) --- není jen logika; produkt agenta, který \emph{má tělo}
a~fyzicky interaguje se světem.

\textbf{Emergence} --- inteligentní chování \emph{vzniká} interakcí mnoha jednoduchých chování;
nikde explicitně naprogramováno.
\end{defbox}

\textbf{Charakteristiky reaktivních agentů:} \emph{behaviorální} (vývoj a~kombinace
jednotlivých chování), \emph{situovaní}, \emph{reaktivní} (jen reagují, nedeliberují).
Reprezentace \textbf{subsymbolická}: konekcionistické sítě, konečné automaty, pravidla IF--THEN.
\subsection{Symbolické reaktivní plánování: subsumpční architektura}
\label{ssec:subsumpce}

\begin{defbox}[Subsumpční architektura (\emph{subsumption architecture})]
Rodney Brooks, 1986--1991; nejúspěšnější reaktivní přístup. Tři Brooksovy teze:
\begin{enumerate}
\item Inteligentní chování \emph{bez symbolické reprezentace}.
\item Inteligentní chování \emph{bez explicitního abstraktního uvažování}.
\item Inteligence = \emph{emergentní vlastnost} určitých složitých systémů.
\end{enumerate}
Inteligence agenta = množina jednoduchých cílově orientovaných \textbf{chování}
(\emph{behaviours}):
\begin{itemize}
\item každé chování = mechanismus výběru akce (dvojice \emph{situace $\to$ akce}): vjemy
  $\to$ akce, jeden dílčí cíl, struktura jednoduchého pravidla;
\item běží \emph{paralelně}, \emph{soupeří} o~kontrolu nad agentem;
\item uspořádána do \textbf{subsumpční hierarchie} = priority (vyšší vrstva \uv{pohlcuje}
  = potlačuje nižší).
\end{itemize}
\end{defbox}

\noindent\textbf{Mechanismus výběru akce v~subsumpční architektuře:}
\begin{algorithmic}[1]
\Function{action}{vjem $p$}
  \State $\mathit{fired} \gets \{\, b \in \mathit{Behaviours} \mid p \text{ splňuje podmínku } b \,\}$
    \Comment{chování, která \uv{vystřelí}}
  \ForAll{$b \in \mathit{fired}$}
    \If{neexistuje $b' \in \mathit{fired}$ takové, že $b' \prec b$}
      \Comment{$b$ není nikým inhibováno}
      \State \Return akce příslušná chování $b$
    \EndIf
  \EndFor
  \State \Return \textbf{null} \Comment{nevystřelilo nic, agent nedělá nic}
\EndFunction
\end{algorithmic}

Formálně: chování = dvojice $\langle c, a\rangle$, $c \subseteq \mathit{Per}$ podmínka,
$a \in A$ akce; na množině chování $R$ relace inhibice $\prec\ \subseteq R \times R$
(striktní totální uspořádání). \textbf{Pozor na směr:} u~Wooldridge i~na přednášce
$b_1 \prec b_2$ = \textbf{$b_1$ inhibuje $b_2$} --- $b_1$ je \emph{níže} v~hierarchii a~má
\emph{přednost}; zápis $b_1 \prec b_2 \prec \dots \prec b_n$ = od nejvyšší priority
k~nejnižší.

\textbf{Vlastnosti:} \emph{jednoduchý} (při desítkách chování záludné priority),
\emph{velmi rychlý} --- hardwarová implementace, konstantní složitost.

\medskip
\noindent\textbf{Příklad: Steelsův průzkumník Marsu.}
Úloha: roj robotů sbírá vzácné vzorky hornin na cizí planetě. Prostředky: navigační signál
základny (robot detekuje jen \emph{gradient}); radioaktivní drobky = \emph{nepřímá
komunikace} (stigmergie); bez přímé komunikace.

\emph{První iterace} --- náhodná procházka a~návrat jednoho robota (priorita klesá shora dolů):
\begin{center}\small
\begin{tabular}{ll}
\toprule
R1 & \textbf{if} vidím překážku \textbf{then} změň směr \\
R2 & \textbf{if} nesu vzorek \textbf{and} jsem na základně \textbf{then} polož vzorek \\
R3 & \textbf{if} nesu vzorek \textbf{and} nejsem na základně \textbf{then} jdi po gradientu nahoru \\
R4 & \textbf{if} vidím vzorek \textbf{then} seber ho \\
R5 & \textbf{if} true \textbf{then} pohybuj se náhodně \\
\bottomrule
\end{tabular}
\end{center}

\emph{Druhá iterace} --- průzkum se stigmergií: vzorky bývají ve shlucích; nálezce cestou domů
rozhazuje drobky = \uv{řekne} to ostatním:
\begin{center}\small
\begin{tabular}{ll}
\toprule
R6 & \textbf{if} nesu vzorek \textbf{and} jsem na základně \textbf{then} polož vzorek \\
R7 & \textbf{if} nesu vzorek \textbf{and} nejsem na základně \textbf{then} polož 2 drobky
     a~jdi po gradientu nahoru \\
R8 & \textbf{if} cítím drobek \textbf{then} seber 1~drobek a~jdi po gradientu dolů \\
\bottomrule
\end{tabular}
\end{center}
Priority (\uv{levé inhibuje pravé}, od nejvyšší):
$\mathrm{R1} \prec \mathrm{R6} \prec \mathrm{R7} \prec \mathrm{R4} \prec
\mathrm{R8} \prec \mathrm{R5}$. Robot na stopě bere po jednom drobku $\to$ stopa se po
vytěžení shluku \emph{vypaří}. Kooperace roje bez modelu světa a~bez zpráv.

\subsection{Konekcionistické reaktivní plánování: agent network architecture}
\label{ssec:maes}

Subsumpce: konflikty řeší \emph{pevná} priorita (rigidní). Konekcionistický přístup:
\textbf{spojitá aktivace šířená sítí} --- výběr akce = dynamika sítě, ne pevné pořadí.
\uv{Reaktivní plánování}: agent se chová, jako by plánoval, explicitní plán nikdy nevytvoří.

\begin{defbox}[Agent network architecture / \emph{spreading activation networks}]
Pattie Maes, 1989--1991. Agent = množina \textbf{kompetenčních modulů} (\emph{competence
modules}), obdoba Brooksových chování. Modul $i$ má:
\begin{itemize}
\item \textbf{předpoklady} $c_i$ (\emph{preconditions}) --- podmínka \emph{proveditelnosti}
  (\emph{executable});
\item \textbf{efekty}: \emph{add list} $a_i$, \emph{delete list} $d_i$ (jako u~STRIPS);
\item \textbf{úroveň aktivace} $\alpha_i$; globální \textbf{práh aktivace} $\theta$.
\end{itemize}
Moduly = orientovaný graf: hrany ze shody pre- a~postpodmínek + hrany časové následnosti
a~konfliktů. Každý krok: šíření aktivace; \textbf{vybrán proveditelný modul s~nejvyšší
aktivací} nad prahem $\theta$.
\end{defbox}

\textbf{Tři zdroje aktivace:}
\begin{itemize}
\item \textbf{stav světa} (\emph{state}) --- moduly se splněnými předpoklady
  (\uv{teď se dá udělat tohle});
\item \textbf{cíle} (\emph{goals}) --- moduly s~cílem v~add listu (\uv{tohle vede k~cíli});
\item \textbf{ochranné cíle} (\emph{protected goals}) --- moduly se splněným cílem v~delete
  listu jsou \emph{inhibovány} (\uv{tohle by zbořilo, co už mám}).
\end{itemize}
\textbf{Tři typy spojů mezi moduly:}
\begin{itemize}
\item \textbf{následnické} (\emph{successor links}): $x \to y$, pokud $y$ přidává předpoklad
  $x$ --- \emph{dopředné} řetězení (\uv{až budu hotov, půjde to dál});
\item \textbf{předchůdcovské} (\emph{predecessor links}): neproveditelný $x$ aktivuje moduly,
  které mu splní chybějící předpoklad --- \emph{zpětné} řetězení (\uv{pomoz mi, ať můžu běžet});
\item \textbf{konfliktní} (\emph{conflictor links}): $x$ \emph{inhibuje} $y$, které ruší jeho
  předpoklady.
\end{itemize}
\textbf{Pět globálních parametrů:}
\begin{itemize}
\item $\theta$ --- práh aktivace (jediný v~systému) = \textbf{uvážlivost vs.\ rychlost}
  (vysoký práh: déle váhá, lépe rozhoduje);
\item $\phi$ --- aktivace vpravená stavem světa; $\gamma$ --- aktivace vpravená cíli;
  poměr $\gamma/\phi$ = \textbf{proaktivita vs.\ reaktivita};
\item $\delta$ --- aktivace odebraná chráněnými cíli;
\item $\pi$ --- \emph{průměrná} aktivace, na kterou se síť po každém kroku normalizuje
  (celková energie konstantní).
\end{itemize}
Žádný proveditelný modul nad prahem $\to$ $\theta$ klesne o~10~\%, znovu --- síť nemůže
\uv{zamrznout}.

\textbf{Jeden krok výpočtu aktivace:} (1)~každý pravdivý fakt rozdělí $\phi$ modulům, které
ho mají v~předpokladech; (2)~každý cíl rozdělí $\gamma$ modulům, které ho přidávají;
chráněné cíle \emph{odeberou} $\delta$ modulům, které by je rušily; (3)~proveditelné moduly
posílají aktivaci následníkům, neproveditelné zpět modulům plnícím chybějící předpoklad,
konfliktní spoje odečítají; (4)~přeškálování na průměr $\pi$. Pak výběr proveditelného
modulu nad $\theta$; provedený modul aktivaci vynuluje.

\emph{Poznámka k~přednášce:} slidy uvádějí práh jako vlastnost modulu (\uv{priorita});
v~Maesově původním modelu je $\theta$ \emph{globální} a~\uv{prioritou} je dynamicky se měnící
aktivace $\alpha_i$.

\medskip\noindent
\textbf{Srovnání se subsumpcí:} pořadí chování ne pevné, ale \emph{dynamické} ze stavu světa
a~cílů. Síť generuje i~vícekrokové \uv{plánovité} sekvence (aktivace řetězená zpět od cíle
předchůdcovskými spoji) bez explicitní datové struktury plánu --- odtud \emph{reaktivní
plánování}. Cena: citlivost na parametry, bez záruky optimality a~úplnosti.

\subsection{Omezení reaktivních architektur}

\begin{itemize}
\item \textbf{Bez modelu světa}: vše z~aktuálních vjemů; paměť nutno \emph{externalizovat}
  změnou prostředí (radioaktivní drobky) --- stigmergie.
\item \textbf{Krátkodobý pohled}: jen aktuální stav a~lokální informace; globální podmínky
  a~dlouhodobé cíle obtížné.
\item \textbf{Emergence není inženýrský postup}: chování se ladí experimentálně, chybí
  metodologie i~verifikace.
\item \textbf{Škálování}: několik vrstev zvládnutelné, \textbf{desítky chování neřešitelné}
  (interakce rostou kvadraticky).
\end{itemize}

\subsection{Hybridní architektury}
\label{ssec:hybridni}

Čistě reaktivní ani čistě deliberativní architektura není ideální $\to$ kombinace:
\begin{itemize}
\item \textbf{deliberativní / plánovací komponenta} --- symbolická úroveň, reprezentace, plány;
\item \textbf{reaktivní komponenta} --- okamžité reakce bez složitých výpočtů.
\end{itemize}
Uspořádání do \textbf{vrstev} (\emph{layered architectures}); reaktivní vrstvy mají zpravidla
\emph{přednost} (bezpečnost před optimalitou).

\begin{defbox}[Horizontální vs.\ vertikální vrstvení]
\textbf{Horizontální vrstvení:} všechny vrstvy \emph{paralelně} připojeny k~senzorům
i~efektorům; každá vrstva = samostatný agent navrhující akci.
\begin{itemize}
\item[$+$] Koncepčně jednoduché: $n$ druhů chování = $n$ vrstev.
\item[$-$] Návrhy vrstev si odporují $\to$ nutná \textbf{mediační funkce} (\emph{mediator
  function}); úzké hrdlo: $m^n$ kombinací ($n$ vrstev, $m$ návrhů), kritický bod selhání.
\end{itemize}

\textbf{Vertikální vrstvení:} vrstvy zapojeny \emph{sériově}.
\begin{itemize}
\item \emph{One-pass}: vjem do nejnižší vrstvy, průchod nahoru, akce z~nejvyšší; přirozená
  hierarchie chování.
\item \emph{Two-pass}: vjemy \emph{zdola nahoru}, řídicí povely (akce) \emph{shora dolů};
  jako řízení ve firmě.
\item[$+$] Bez mediace, lineární složitost interakcí ($n-1$ rozhraní).
\item[$-$] Křehké: selhání kterékoli vrstvy shodí celého agenta.
\end{itemize}
\end{defbox}

\begin{defbox}[TouringMachines (Ferguson, 1992)]
\textbf{Horizontální} třívrstvá architektura; simulovaný agent v~dopravním prostředí.
\begin{itemize}
\item \textbf{Modelovací vrstva} (\emph{modelling layer}) --- symbolické modely agenta,
  prostředí a~ostatních agentů; předvídá a~řeší konflikty, stanovuje cíle pro plánovací vrstvu.
\item \textbf{Plánovací vrstva} (\emph{planning layer}) --- proaktivní chování; výběr
  z~předpřipravených plánů (jako PRS), sestavení plánu cesty.
\item \textbf{Reaktivní vrstva} (\emph{reactive layer}) --- pravidla \emph{situace $\to$ akce},
  okamžitá reakce (vyhýbání překážkám).
\end{itemize}
Vrstvy běží paralelně; výstupy sjednocují \textbf{řídicí pravidla} (\emph{control rules})
= konkrétní mediační funkce; potlačují vstupy a~výstupy vrstev (\emph{censor}
a~\emph{suppressor} pravidla).
\end{defbox}

\begin{defbox}[InteRRaP (Müller, 1994)]
\textbf{Vertikální dvouprůchodová} (\emph{two-pass}) třívrstvá architektura.
\begin{itemize}
\item \textbf{Vrstva kooperativního plánování} (\emph{cooperative planning layer}) ---
  sociální interakce, společné plány s~ostatními agenty.
\item \textbf{Vrstva lokálního plánování} (\emph{local planning layer}) --- individuální
  plánování.
\item \textbf{Behaviorální vrstva} (\emph{behaviour-based layer}) --- reaktivní chování.
\end{itemize}
Každá vrstva má \emph{vlastní znalostní bázi} své úrovně abstrakce (svět / plány / společné
plány a~modely ostatních agentů). Dvě operace mezi vrstvami: \emph{aktivace zdola nahoru}
(nižší vrstva situaci nezvládne, předá výš), \emph{zásah shora dolů} (vyšší vrstva vykonává
svá rozhodnutí přes nižší).
\end{defbox}

\medskip
\noindent\textbf{Reálný příklad: Stanley} (Stanford) --- autonomní vůz, vítěz DARPA Grand
Challenge 2005. Vrstvy: senzory $\to$ abstraktní percepce $\to$ plánování a~řízení
$\to$ rozhraní vozidla; navíc vrstvy uživatelského rozhraní a~globálních služeb.

\subsection{Komplexní architektury: IDA a~globální pracovní prostor}
\label{ssec:ida}

Opak reaktivního minimalismu: integrovat do jednoho agenta \emph{všechno} --- výběr akce,
paměť, deliberaci, emoce.

\begin{defbox}[IDA --- \emph{Intelligent Distribution Agent} (S.~Franklin)]
Jedna z~nejsložitějších postavených agentních architektur. Reálná úloha amerického
námořnictva: přidělování personálu na nová služební místa (odtud \uv{distribution});
distribuovaná je i~architektura sama.

\textbf{Východisko: teorie globálního pracovního prostoru} (\emph{global workspace theory},
Baars, 1988--97) --- jedna z~teorií vědomí:
\begin{itemize}
\item mysl \emph{je} MAS: mnoho jednoduchých, převážně nevědomých specializovaných procesů;
\item komunikace \emph{zřídka}, jen přes sdílenou paměť typu \emph{blackboard} organizovanou
  jako asociativní pole (srov.\ kap.~\ref{sec:kooperace});
\item procesy dynamicky tvoří \textbf{koalice} vyšších řádů;
\item \textbf{výběr akce}: koalice nejlépe odpovídající aktuálním vstupům se dostane do
  \uv{vědomí} (globálního pracovního prostoru) a~vykoná se;
\item \emph{hierarchie kontextů} --- kontext cílů, vjemů, smyslů, konceptů, kulturní kontext.
\end{itemize}
\textbf{Hodnocení:} kombinuje mnoho přístupů z~MAS i~zbytku UI (výběr akce, paměť,
deliberace, emoce) --- zároveň slabina: řada rozhodnutí \emph{ad hoc}, sladění heterogenních
modelů do optimálně spolupracujícího celku velmi obtížné.
\end{defbox}

\slajd{mas-ida.png}{0.96}{Architektura IDA. Modré bloky jsou moduly vnímání a~paměti, zelené
deliberace a~vyjednávání, červené učení a~výstup; vše se potkává v~\uv{Consciousness}
(globálním pracovním prostoru) a~v~\emph{Behavior Net}, který vybírá akci.}

IDA = \emph{čtvrtý} typ arbitráže mezi chováními: subsumpce \textbf{pevná priorita}, síť Maes
\textbf{úroveň aktivace}, vrstvené architektury \textbf{mediační funkce}, IDA \textbf{soutěž
dynamicky vznikajících koalic} o~přístup do globálního pracovního prostoru.

\subsection{Srovnání architektur}

\begin{center}
\small
\begin{tabular}{@{}p{0.15\textwidth}
  >{\raggedright\arraybackslash}p{0.19\textwidth}
  >{\raggedright\arraybackslash}p{0.19\textwidth}
  >{\raggedright\arraybackslash}p{0.17\textwidth}
  >{\raggedright\arraybackslash}p{0.17\textwidth}@{}}
\toprule
& \textbf{Deduktivní} & \textbf{BDI / PRS} & \textbf{Reaktivní} & \textbf{Hybridní} \\
\midrule
Reprezentace & symbolická (FOL) & symbolická (beliefs, plány) & subsymbolická / pravidla
  & smíšená po vrstvách \\
Výběr akce & dokazování $\mathit{Do}(a)$ & deliberace + plán z~knihovny
  & priorita / aktivace & vrstvy + arbitráž \\
Model světa & úplný & částečný, revidovatelný & žádný & po vrstvách \\
Rychlost & velmi nízká & střední & vysoká (konst.) & vysoká v~reflexu \\
Reaktivita & špatná & střední (laditelná) & výborná & dobrá \\
Proaktivita & dobrá & výborná & žádná & dobrá \\
Návrh & obtížný, ale formální & systematický & experimentální & složitý \\
\bottomrule
\end{tabular}
\end{center}

\subsection{Shrnutí ke~zkoušce}

\begin{itemize}
\item \textbf{Deduktivní agent}: $\mathit{DB} \in 2^L$, $\mathit{DB} \vdash_\rho \mathit{Do}(a)$;
  dvoufázový výběr akce (dokazatelná akce $\to$ konzistentní akce $\to$ nic).
  Elegantní, nepraktický; problém transdukce a~vypočitatelné racionality.
\item \textbf{BDI}: dvě fáze --- deliberace ($\mathit{brf}$, $\mathit{options}$, $\mathit{filter}$)
  a~means-ends reasoning (plánování). Beliefs subjektivní a~omylné, desires mohou být
  nekonzistentní, intentions = závazky: přetrvávají, omezují další deliberaci.
\item \textbf{STRIPS}: deskriptor akce $\langle P_a,D_a,A_a\rangle$, plánovací problém
  $\langle B_0,O,G\rangle$, $B_i = (B_{i-1}\setminus D_{a_i})\cup A_{a_i}$;
  plán \emph{přípustný} ($B_{i-1}\models P_{a_i}$) vs.\ \emph{korektní} (navíc $B_n \models G$).
  Umět projít svět kostek.
\item \textbf{BDI smyčka} a~role $\mathit{reconsider}$ a~$\mathit{sound}$.
\item \textbf{Commitment}: blind / single-minded / open-minded; bold vs.\ cautious agent;
  efektivita $=$ dosažené/všechny záměry; pomalé prostředí $\to$ bold, rychlé $\to$ cautious.
\item \textbf{PRS}: první implementace BDI; plán = Goal + Context + Body; knihovna plánů místo
  plánování; zásobník záměrů; deliberace metaplány nebo utilitou; potomci AgentSpeak/Jason,
  JAM, JACK, Jadex.
\item Reaktivní přístup: trojice \emph{situatedness --- embodiment --- emergence};
  odmítá symbolickou reprezentaci.
\item \textbf{Subsumpce} (Brooks) = \emph{symbolické} reaktivní plánování: jednoduchá chování
  \emph{situace~$\to$~akce}, paralelní, pevná priorita.
  Umět Steelsova průzkumníka Marsu a~roli stigmergie (drobky).
\item \textbf{Agent network architecture} (Maes) = \emph{konekcionistické} reaktivní
  plánování: moduly s~předpoklady a~efekty (add, delete); aktivace se šíří následnickými,
  předchůdcovskými a~konfliktními spoji; zdroje: stav světa, cíle, inhibice od chráněných
  cílů; vybere se proveditelný modul s~nejvyšší aktivací nad prahem. $\gamma/\phi$ ladí
  proaktivitu vs.\ reaktivitu.
\item Omezení reaktivních agentů: bez modelu světa, krátkodobý pohled, neinženýrský návrh,
  neškálovatelné množství vrstev.
\item \textbf{Hybridní}: horizontální (paralelní vrstvy + mediátor, $m^n$ kombinací, úzké hrdlo)
  vs.\ vertikální (sériové, one-pass/two-pass, $n-1$ rozhraní, křehké).
\item \textbf{TouringMachines} = horizontální 3~vrstvy (modelling, planning, reactive) + řídicí
  pravidla. \textbf{InteRRaP} = vertikální two-pass 3~vrstvy (kooperativní plánování, lokální
  plánování, behaviorální) + vlastní KB na vrstvu. \textbf{Stanley} = reálné nasazení.
\item \textbf{IDA} (Franklin) + \emph{global workspace theory} (Baars): mysl jako MAS
  jednoduchých nevědomých procesů komunikujících přes blackboard; koalice soutěží o~vstup do
  \uv{vědomí}, vítězná se vykoná; hierarchie kontextů. Nejsložitější architektura, mnoho
  ad hoc rozhodnutí.
\item Čtyři typy arbitráže mezi chováními: \textbf{pevná priorita} (subsumpce),
  \textbf{úroveň aktivace} (Maes), \textbf{mediační funkce} (vrstvené), \textbf{soutěž koalic}
  (IDA).
\end{itemize}
\section{Problém hledání cesty}
\label{sec:cesta}
%=====================================================================

\noindent\mimoq{celá kapitola} NAIL106 plánování i~robotiku výslovně \emph{nepokrývá}, slidy
k~hledání cesty nemají nic; okruh ho ale jmenuje $\Rightarrow$ nutné umět. Minimum, které na
zkoušce obstojí: A* s~vlastnostmi heuristiky + \emph{multiagentní} část --- dynamické prostředí,
koordinace více agentů.

\subsection{Formulace a~zařazení}

Hledání cesty (\emph{pathfinding}) --- nejčastější konkrétní podúloha situovaného agenta:
\begin{itemize}
\item kde: součást plánovací vrstvy hybridní architektury, akce \texttt{move\_towards(P)}
  v~BDI plánu, nebo služba poskytovaná jiným agentům;
\item multiagentní specifika: prostředí \textbf{dynamické} (mapa se mění), ostatní agenti =
  \textbf{pohyblivé překážky}, nutná koordinace.
\end{itemize}

\begin{defbox}[Problém hledání cesty]
Ohodnocený graf $G = (V,E)$, nezáporné ohodnocení hran $w : E \to \mathbb{R}_0^+$, počáteční
vrchol $s \in V$, cílový vrchol $g \in V$. Hledáme cestu $s = v_0, v_1, \dots, v_k = g$
minimalizující $\sum_{i} w(v_{i-1},v_i)$.

Graf bývá zadán \emph{implicitně} nástupnickou funkcí (mřížka 4-/8-souvislá, navigační síť,
graf viditelnosti).
\end{defbox}

\mimo{} \textbf{Metaheuristiky} --- \emph{mravenčí algoritmy} (ACO) a~\emph{optimalizace rojem
částic} (PSO) z~okruhu \emph{Přírodou inspirované počítání}:
\begin{itemize}
\item ACO: inspirováno hledáním nejkratší cesty k~potravě; směrování v~sítích (AntNet),
  cesta robota v~mřížce;
\item PSO: v~robotice cesta = vektor průjezdových bodů, spojitě optimalizována délka,
  hladkost, odstup od překážek;
\item \textbf{pro zadání výše se nevyplatí}: statický graf se známými cenami $\Rightarrow$ A*
  \emph{optimální} a~rychlejší, metaheuristika bez záruky optimality;
\item smysl až u~úloh těžších než nejkratší cesta: spojitý prostor s~překážkami, více kritérií,
  silně dynamické či stochastické ceny, směrování bez globální znalosti topologie.
\end{itemize}

\subsection{A*}

Prohledávání typu \emph{best-first}; pro každý vrchol $n$:
\begin{align*}
g(n) &= \text{cena nejlepší dosud nalezené cesty } s \to n, \\
h(n) &= \text{heuristický odhad ceny } n \to g_{\mathrm{cil}}, \\
f(n) &= g(n) + h(n) \quad \text{--- odhad ceny nejlepší cesty } s \to g_{\mathrm{cil}}
  \text{ vedoucí přes } n .
\end{align*}

\noindent\textbf{Algoritmus A*} (Hart, Nilsson, Raphael, 1968):
\begin{algorithmic}[1]
\State $\mathit{OPEN} \gets \{s\}$ (prioritní fronta podle $f$), $\mathit{CLOSED} \gets \emptyset$
\State $g(s) \gets 0$, $f(s) \gets h(s)$, ostatní $g \gets \infty$
\While{$\mathit{OPEN} \ne \emptyset$}
  \State $n \gets \arg\min_{x \in \mathit{OPEN}} f(x)$; odeber $n$ z~$\mathit{OPEN}$
  \If{$n$ je cíl} \Return rekonstruovaná cesta \EndIf
  \State přidej $n$ do $\mathit{CLOSED}$
  \ForAll{$m \in \mathit{succ}(n)$}
    \State $g_{\mathrm{new}} \gets g(n) + w(n,m)$
    \If{$g_{\mathrm{new}} < g(m)$}
      \State $g(m) \gets g_{\mathrm{new}}$; $\mathit{parent}(m) \gets n$;
             $f(m) \gets g(m) + h(m)$
      \State vlož/aktualizuj $m$ v~$\mathit{OPEN}$ (a~odeber z~$\mathit{CLOSED}$, byl-li tam)
    \EndIf
  \EndFor
\EndWhile
\State \Return \uv{cesta neexistuje}
\end{algorithmic}

\begin{defbox}[Přípustnost a~konzistence heuristiky]
\textbf{Přípustná} (\emph{admissible}): nikdy nenadhodnocuje, $h(n) \le h^*(n)$ pro všechna
$n$; $h^*(n)$ = skutečná cena nejlevnější cesty z~$n$ do cíle.

\textbf{Konzistentní} (\emph{consistent}, monotónní): trojúhelníková nerovnost pro každou
hranu $(n,m)$
\[
h(n) \le w(n,m) + h(m), \qquad h(g_{\mathrm{cil}}) = 0 .
\]
\textbf{Vlastnosti:} konzistence $\Rightarrow$ přípustnost. Přípustná $h$ $\Rightarrow$ A*
\emph{optimální} (najde nejlevnější cestu). Konzistentní $h$ $\Rightarrow$ navíc $f$ podél
každé cesty neklesající a~každý vrchol expandován \emph{nejvýše jednou}.
\end{defbox}

\textbf{Speciální případy a~varianty:}
\begin{itemize}
\item $h \equiv 0$ --- A* degeneruje na \textbf{Dijkstrův algoritmus};
\item heuristiky na mřížce: manhattanská $|\Delta x| + |\Delta y|$ (4~směry), oktilová
  (8~směrů), eukleidovská;
\item \textbf{Weighted A*}: $f = g + \varepsilon h$, $\varepsilon > 1$ --- rychlejší, cesta
  nejvýše $\varepsilon\times$ dražší;
\item \textbf{IDA*}: šetří paměť ($O(d)$ místo exponenciální).
\end{itemize}

\subsection{Dynamické prostředí: D*~Lite}

Agent za pohybu zjišťuje změny mapy (nová překážka, změněná cena). A* od nuly v~každém kroku
= drahé $\Rightarrow$ \textbf{inkrementální} algoritmy: při lokální změně přepočítají jen
zasaženou část.

\begin{defbox}[LPA* a~D*~Lite]
\textbf{LPA*} (\emph{Lifelong Planning A*}, Koenig, Likhachev, 2001) --- inkrementální A* pro
pevný start a~cíl. Ke každému vrcholu kromě $g(n)$ i~\emph{lookahead} hodnota
\[
\mathit{rhs}(n) = \begin{cases}
0 & n = s,\\
\min_{p \in \mathit{pred}(n)} \big(g(p) + w(p,n)\big) & \text{jinak.}
\end{cases}
\]
Vrchol \textbf{lokálně konzistentní} $\Leftrightarrow$ $g(n) = \mathit{rhs}(n)$. Změna ceny
hrany $\Rightarrow$ některé vrcholy nekonzistentní $\Rightarrow$ do prioritní fronty; změna se
propaguje jen tam, kde je potřeba.

\textbf{D*~Lite} (Koenig, Likhachev, 2002) --- LPA* prohledávaný \emph{od cíle ke startu}
(hodnoty zůstanou použitelné, když se robot posune) + korekce prioritní fronty při pohybu.
Funkčně ekvivalentní staršímu \textbf{D*} (Stentz, 1994), ale výrazně jednodušší; \emph{de
facto} standard v~mobilní robotice.
\end{defbox}

Neznámé prostředí: \emph{optimistický předpoklad} (\emph{freespace assumption}) --- neznámé
buňky = průchodné, agent jde po plánu, po nárazu přeplánuje; přesně scénář, pro nějž byl D*
navržen.

\subsection{Hledání cest pro více agentů (MAPF)}
\label{ssec:mapf}

\noindent\mimoq{není ve slidech, ale je to jediná vysloveně multiagentní část okruhu}

\begin{defbox}[Multi-Agent Path Finding, MAPF]
Vstup: graf $G=(V,E)$ a~$k$ agentů, agent $i$ má start $s_i$ a~cíl $g_i$. Čas diskrétní;
v~každém kroku každý agent buď přejde po hraně, nebo počká na místě.

\textbf{Konflikty} (řešení musí být \emph{bezkonfliktní}): \emph{vrcholový konflikt} --- dva
agenti ve stejném čase ve stejném vrcholu; \emph{hranový / \uv{swap} konflikt} --- dva agenti
si v~jednom kroku vymění pozice.

\textbf{Účelová funkce:} \emph{sum-of-costs} $\sum_i \mathit{cost}_i$, nebo \emph{makespan}
$\max_i \mathit{cost}_i$. Nalezení \emph{optimálního} řešení: NP-těžké.
\end{defbox}

\textbf{Prioritizované plánování} (rozvázaný přístup), varianta \emph{Cooperative A*}
(Silver, 2005):
\begin{itemize}
\item agenti uspořádáni podle priority; pro $i = 1,\dots,k$ cesta agenta $i$ pomocí A*
  v~\emph{prostoru--čase} (stav = dvojice $\langle$vrchol, čas$\rangle$), s~vyhnutím již
  naplánovaným trajektoriím v~\textbf{rezervační tabulce};
\item rychlé, škáluje, ale \textbf{neúplné} --- existují řešitelné instance, které neřeší žádné
  pořadí priorit (dva agenti v~úzkém koridoru, jeden musí uhnout do výklenku).
\end{itemize}

\begin{defbox}[Conflict-Based Search (CBS)]
Sharon a~kol., 2012--2015. Dvouúrovňový, \emph{optimální} a~\emph{úplný}; vyhýbá se
exponenciálnímu sdruženému stavovému prostoru $V^k$.

\textbf{Nízká úroveň:} jeden agent, optimální cesta v~prostoru--čase pomocí A* při
respektování množiny \emph{omezení} tvaru $\langle a_i, v, t\rangle$ (\uv{agent $i$ nesmí být
ve vrcholu $v$ v~čase $t$}).

\textbf{Vysoká úroveň:} \emph{strom omezení} prohledávaný podle ceny. Kořen = prázdná množina
omezení + individuálně optimální cesty. Vezmi uzel s~nejnižší cenou: cesty bezkonfliktní
$\Rightarrow$ řešení optimální; jinak první konflikt $\Rightarrow$ \emph{rozvětvení} na dva
potomky, do každého omezení pro jednoho z~kolidujících agentů a~přeplánuje se \emph{jen} on.
\end{defbox}

\subsection{Shrnutí ke~zkoušce}

\begin{itemize}
\item A*: $f = g+h$; \emph{přípustnost} $h \le h^*$ $\Rightarrow$ optimalita;
  \emph{konzistence} $h(n) \le w(n,m)+h(m)$ $\Rightarrow$ každý vrchol se expanduje jednou.
  $h\equiv 0$ dá Dijkstru; weighted A* je $\varepsilon$-suboptimální.
\item Dynamické prostředí: LPA* (hodnoty $g$ a~\emph{rhs}, lokální konzistence
  $g=\mathit{rhs}$), D*~Lite (prohledávání od cíle, inkrementální oprava), freespace assumption.
\item MAPF: vrcholové a~hranové (swap) konflikty, sum-of-costs vs.\ makespan;
  optimální MAPF je NP-těžký.
\item Prioritizované plánování / Cooperative A* s~rezervační tabulkou --- rychlé, ale neúplné.
  \textbf{CBS} --- nízká úroveň A* s~omezeními, vysoká úroveň větvení stromu omezení podle
  konfliktu; optimální a~úplný.
\end{itemize}

%=====================================================================
\section{Komunikace a~znalosti v~MAS: ontologie, řečové akty, FIPA-ACL, protokoly}
\label{sec:komunikace}
%=====================================================================
\subsection{Od agenta k~multiagentnímu systému}

Jeden agent umíme; systém agentů vyžaduje komunikaci, a~tedy:
\begin{itemize}
\item \textbf{reprezentaci znalostí} --- v~čem jsou znalosti vyjádřeny;
\item \textbf{společný slovník} --- stejné slovo = stejný význam (\emph{ontologie});
\item \textbf{standardní zprávy} --- syntaxe a~sémantika komunikace (\emph{ACL});
\item \textbf{předvídatelné chování při komunikaci} --- \emph{interakční protokoly};
\item \textbf{technické problémy} --- distribuovanost, mobilita agentů, asynchronnost,
  nespolehlivé kanály.
\end{itemize}

\subsection{Co je ontologie}

\begin{defbox}[Ontologie]
\emph{Filozoficky}: nauka o~bytí, Aristotelova \uv{první filozofie}.

\emph{V~informatice}: \textbf{explicitní a~formalizovaný popis části reality} --- formální
deklarativní popis s~\emph{glosářem} (definice konceptů) a~\emph{tezaurem} (definice vztahů
mezi koncepty); slovník pro standardní ukládání a~výměnu znalostí o~doméně.

\textbf{Gruber (1995):} \uv{Ontologie je popis (podobný formální specifikaci programu) konceptů
a~vztahů, které mohou pro agenta nebo komunitu agentů formálně existovat.}
Zkráceně: \emph{ontologie je explicitní specifikace konceptualizace}.
\end{defbox}

\textbf{Motivační příklad:} věta \uv{7777 je nové CD kapely SRPR} vyžaduje sdílenou
konceptualizaci: 7777 = album (žánry, skladby, interpret); SRPR = kapela (členové, nástroje\dots).

\begin{defbox}[Formalismus ontologie]
\begin{itemize}
\item \textbf{Třídy} (\emph{classes}, koncepty) --- věci s~něčím společným.
\item \textbf{Instance} (\emph{individuals}, objekty) --- konkrétní jedinci tříd.
\item \textbf{Vlastnosti} (\emph{properties}) --- atributy tříd a~instancí.
\item \textbf{Relace mezi třídami}, zejména \textbf{podtřída} (\emph{is-a}) --- \emph{tranzitivní}.
\item Další relace a~vlastnosti (základní znalosti o~doméně).
\end{itemize}
\textbf{Strukturální část} (třídy, relace, axiomy) = vlastní \emph{ontologie};
+ \textbf{fakta o~konkrétních jedincích} = \textbf{znalostní báze} (\emph{knowledge base}).
\end{defbox}

\textbf{Ontologie ontologií podle síly formalizace} (od slabých k~silným):
slovník (řízený slovník termínů) $\to$ glosář (významy v~přirozeném jazyce)
$\to$ tezaurus (synonyma) $\to$ neformální hierarchie (Amazon, wiki) $\to$
formální \emph{is-a} hierarchie (subsumpce tříd) $\to$ třídy s~vlastnostmi $\to$
omezení hodnot (\uv{každý člověk má právě jednu matku}) $\to$ libovolná logická omezení
(složité odvozovací algoritmy).

\textbf{Ontologie ontologií podle rozsahu:}
\begin{itemize}
\item \textbf{Aplikační} --- v~praxi nejběžnější, těžko znovupoužitelné.
\item \textbf{Doménové} --- předmět výzkumu sémantického webu.
\item \textbf{Horní} (\emph{upper ontologies}) --- popis \uv{všeho} (Thing, Living thing, \dots),
  základ doménových ontologií: BFO, GFO, UFO; WordNet a~IDEAS nevhodné pro formální odvozování.
\end{itemize}

\subsection{Ontologické jazyky}

Formální deklarativní jazyky pro reprezentaci znalostí: fakta + odvozovací pravidla; základ
nejčastěji logika prvního řádu nebo \emph{deskripční logika}.

\textbf{Rámce} (\emph{frames}) --- historický předchůdce (Minsky), klasická UI a~expertní
systémy; obdoba objektového programování:

\begin{center}\small
\begin{tabular}{ll}
\toprule
\textbf{Terminologie rámců} & \textbf{Objektová terminologie} \\
\midrule
Frame (rámec) & třída objektů \\
Slot & vlastnost / atribut objektu \\
Trigger, accessor, mutator & metoda \\
\bottomrule
\end{tabular}
\end{center}

\textbf{XML} --- pro ontologie nevznikl, ale používá se: vlastní značky
(\texttt{<product>}, \texttt{<title>}, \texttt{<artist>}) = slovník. Chybí sémantika ---
popisuje jen strukturu dokumentu, ne význam.

\begin{defbox}[RDF --- \emph{Resource Description Framework}]
Standard reprezentace znalostí (nejen) pro web. \textbf{Jednoduchý}: málo expresivní, rychlé
odvozování.

Vše jako \textbf{trojice} \emph{subjekt -- predikát -- objekt},
$\mathit{Predikát}(\mathit{Subjekt},\mathit{Objekt})$; např.\
$\mathit{MarriedTo}(\mathit{Karel},\mathit{Jája})$,
$\mathit{FatherOf}(\mathit{Karel},\mathit{Péťa})$.
Množina trojic = orientovaný ohodnocený graf (\emph{RDF graph}); zdroje identifikovány IRI.
Serializace: RDF/XML, Turtle, N-Triples, JSON-LD. Dotazování: SPARQL.
Nadstavba \textbf{RDFS}: třídy, \texttt{subClassOf}, \texttt{domain}, \texttt{range}.
\end{defbox}

\begin{defbox}[OWL --- \emph{Web Ontology Language}]
Původ v~sémantickém webu; verze 1 a~2. \emph{Kolekce několika formalismů} pro popis
ontologií; syntaxe XML, RDF, funkcionální notace, Manchester syntax nebo Turtle. Cíl:
formální, a~přesto prakticky použitelný jazyk.

Základ: \textbf{deskripční logika} --- \emph{rozhodnutelný fragment} logiky prvního řádu.

\textbf{Předpoklad otevřeného světa} (\emph{open world assumption}, OWA): co nevíme, je
\emph{nedefinované}; oproti uzavřenému světu (databáze, Prolog), kde neznámé = \emph{nepravdivé}.
\end{defbox}

\begin{defbox}[Deskripční logika]
Entity: \textbf{koncepty} (třídy), \textbf{role} (vlastnosti, predikáty),
\textbf{individua} (objekty). \textbf{Axiom} = logický výraz o~rolích a~konceptech; oproti
rámcům a~OOP (třída popsána \emph{úplně}) popisují axiomy třídu jen \emph{částečně}
(mohou přicházet z~více zdrojů).

Rodina systémů podle povolených konstruktorů. Základní $\mathcal{ALC}$: negace, konjunkce,
disjunkce konceptů, omezené existenční a~univerzální kvantifikátory. Rozšíření: hierarchie
rolí ($\mathcal{H}$), tranzitivita ($\mathcal{S}$), inverzní role ($\mathcal{I}$),
kardinalitní omezení ($\mathcal{N}$, $\mathcal{Q}$), nominály ($\mathcal{O}$).

Znalostní báze: \textbf{T-Box} (terminologie --- axiomy o~konceptech a~rolích)
+ \textbf{A-Box} (asertivní část --- fakta o~individuích).
\end{defbox}

\textbf{Profily OWL.} \emph{OWL~1 Lite} ($\mathcal{SHIF}$) --- nejjednodušší, blízko RDF,
omezení kvůli rychlému strojovému zpracování; \emph{OWL~1 DL} ($\mathcal{SHOIN}$) --- odpovídá
deskripční logice (např.\ disjunktnost tříd), úplný a~rozhodnutelný; \emph{OWL~1 Full} ---
největší vyjadřovací síla, mnoho problémů nerozhodnutelných. \emph{OWL~2} ($\mathcal{SROIQ}$)
+ profily: \emph{EL} (polynomiální odvozování; velké lékařské ontologie typu SNOMED),
\emph{QL} (dotazy nad znalostní bází, napojení na relační databáze), \emph{RL} (axiomy ve
tvaru pravidel).

\begin{defbox}[KIF --- \emph{Knowledge Interchange Format}]
Logika prvního řádu, LISPovská prefixová syntaxe. Projekt DARPA \emph{Knowledge Sharing
Effort}: KIF = \textbf{vnitřní jazyk} zpráv (obsah), KQML = \emph{vnější} (obálka). Příklady:
\begin{quote}\small\ttfamily
(salary 015-46-3946 widgets 72000)\\
(forall (?F ?T) (=> (and (instance ?F Farmer) (instance ?T Tractor)) (likes ?F ?T)))
\end{quote}
\end{defbox}

\textbf{DAML+OIL} --- DARPA Agent Markup Language + Ontology Interchange Language, přímý
předchůdce OWL, opuštěn 2006.

\subsection{Uvažování o~znalostech: epistemická logika}

\noindent\mimoq{slidy modální logiky výslovně vynechávají; okruh ale mluví o~\uv{znalostech
v~MAS}, takže minimum se hodí umět}

\begin{defbox}[Epistemická logika a~společná znalost]
Znalost agenta = modální operátor $K_i\varphi$ (\uv{agent $i$ \emph{ví}, že $\varphi$}),
\textbf{kripkovská sémantika možných světů}: $K_i\varphi$ platí ve světě $w$ $\iff$ $\varphi$
platí ve \emph{všech} světech, které $i$ neumí od $w$ odlišit.
\textbf{Znalost}: systém \textbf{S5}, klíčový axiom \textbf{T}: $K_i\varphi \Rightarrow \varphi$
(co vím, platí). \textbf{Přesvědčení} ($B_i$, srov.\ BDI): systém \textbf{KD45}, místo T slabší
\textbf{D}: $\neg B_i\bot$ --- konzistentní, ale \emph{může být mylné}. Znalost vs.\
přesvědčení = přesně axiom~T.

Skupina $G$: $E_G\varphi$ --- každý v~$G$ ví $\varphi$; \textbf{společná znalost}
$C_G\varphi \equiv E_G\varphi \wedge E_G E_G\varphi \wedge \dots$ (všichni vědí, že všichni
vědí, \dots).

\textbf{Koordinovaný útok} (problém dvou generálů): domluva současného útoku přes
nespolehlivého posla. Žádný konečný počet potvrzení nevytvoří společnou znalost --- po $k$
zprávách nejvýše $E_G^k\varphi$, nikdy $C_G\varphi$. Poučení: \emph{nespolehlivá komunikace
nikdy nevytvoří společnou znalost}, tedy ani dokonalou koordinaci.
\end{defbox}
\subsection{Proč je komunikace agentů jiná}

OOP: \uv{komunikace} = volání metody, volaný \emph{musí} vyhovět. Agenti nemohou jiného agenta
přinutit ani mu měnit vnitřní proměnné --- \textbf{provádějí komunikační akt} s~cílem
\begin{itemize}
\item vyměnit si informace, nebo
\item \emph{ovlivnit} jiné agenty, aby něco udělali.
\end{itemize}
Příjemce má vlastní cíle a~s~informací, žádostí či dotazem naloží po svém:
\textbf{komunikace v~MAS je akce, ne volání}.

\subsection{Teorie řečových aktů}

\begin{defbox}[Řečové akty (\emph{speech acts}) --- Austin, 1962]
Části užívání jazyka mají charakter \textbf{akcí} --- mění stav světa jako fyzické akce:
\begin{itemize}
\item \uv{Prohlašuji vás mužem a~ženou.}
\item \uv{Vyhlašuji válku.}
\end{itemize}
\emph{Pragmatická} teorie jazyka: řeč jako prostředek k~dosahování cílů.
Typická performativní slovesa: \emph{request}, \emph{inform}, \emph{promise}.
\end{defbox}

\begin{defbox}[Tři vrstvy řečového aktu]
\begin{itemize}
\item \textbf{Lokuční akt} (\emph{locutionary act}) --- \emph{co bylo řečeno}; samotný výrok.
  \uv{Udělej mi čaj.}
\item \textbf{Ilokuční akt} (\emph{illocutionary act}) --- \emph{co bylo míněno};
  lokuce + performativní význam (dotaz, žádost, \dots). \uv{Požádal mě o~čaj.}
\item \textbf{Perlokuční akt} (\emph{perlocutionary act}) --- \emph{co se skutečně stalo};
  efekt aktu. \uv{Přiměla mě, abych jí uvařil čaj.}
\end{itemize}
Pro MAS klíčová \textbf{ilokuce} --- obsah zprávy (performativa).
\end{defbox}

\textbf{Searle} --- podmínky \uv{zdařilosti} řečového aktu; pro \emph{MLUVČÍ požádá
POSLUCHAČE o~akci}:
\begin{itemize}
\item \textbf{Standardní vstupně-výstupní podmínky}: posluchač slyší, nejde o~scénu ve filmu\dots
\item \textbf{Předpoklady} (aby mluvčí mohl akt zvolit): posluchač je schopen akci provést;
  mluvčí věří, že je schopen; není zřejmé, že by ji provedl i~bez požádání.
\item \textbf{Upřímnost} (\emph{sincerity}): mluvčí opravdu chce, aby se akce provedla.
\end{itemize}

\begin{defbox}[Searlovy kategorie řečových aktů]
Starší dělení: \emph{request}, \emph{advice}, \emph{statement}, \emph{promise}, \dots

Novější (standardní) pět tříd:
\begin{itemize}
\item \textbf{Asertiva} (\emph{assertives}, reprezentativa) --- informují (\emph{inform}).
\item \textbf{Direktiva} (\emph{directives}) --- žádají o~akci (\emph{request}).
\item \textbf{Komisiva} (\emph{commissives}) --- mluvčí se zavazuje (\emph{promise}).
\item \textbf{Expresiva} (\emph{expressives}) --- mentální stav, emoce (\uv{děkuji}).
\item \textbf{Deklarace} (\emph{declarations}) --- mění stav věcí (válka, sňatek).
\end{itemize}
Dvě složky každého aktu: \textbf{performativní sloveso} (request, query, inform)
+ \textbf{propoziční obsah} (\uv{okno je zavřené}).
\end{defbox}

\begin{defbox}[Plánovací teorie řečových aktů (Cohen, Perrault, 1979)]
\uv{\dots modelování [řečových aktů] v~plánovacím systému jako operátorů definovaných\dots{}
pomocí přesvědčení a~cílů mluvčích a~posluchačů. S~řečovými akty se tak zachází stejně jako
s~fyzickými akcemi.}

Sémantika ve stylu \textbf{STRIPS} (předpoklady -- delete -- add) s~modálními operátory pro
\emph{beliefs}, \emph{abilities}, \emph{wants}.
\end{defbox}

\noindent\textbf{Příklad: Request a~Inform.}
\begin{center}\small
\begin{tabular}{p{0.16\textwidth} p{0.38\textwidth} p{0.38\textwidth}}
\toprule
& \textbf{Request(S,H,A)} & \textbf{Inform(S,H,$\varphi$)} \\
\midrule
Předpoklad \emph{cando}
  & $(S\ \mathrm{believe}\ (H\ \mathrm{cando}\ A))\ \wedge$ \newline
    $(S\ \mathrm{believe}\ (H\ \mathrm{believe}\ (H\ \mathrm{cando}\ A)))$
  & $(S\ \mathrm{believe}\ \varphi)$ \\
Předpoklad \emph{want}
  & $(S\ \mathrm{believe}\ (S\ \mathrm{want}\ \mathrm{requestInstance}))$
  & $(S\ \mathrm{believe}\ (S\ \mathrm{want}\ \mathrm{informInstance}))$ \\
Efekt
  & $(H\ \mathrm{believe}\ (S\ \mathrm{believe}\ (S\ \mathrm{want}\ A)))$
  & $(H\ \mathrm{believe}\ (S\ \mathrm{believe}\ \varphi))$ \\
\bottomrule
\end{tabular}
\end{center}
Efekt \emph{Inform} \emph{není} \uv{$H$ uvěří $\varphi$}, jen \uv{$H$ uvěří, že $S$ věří
$\varphi$} --- autonomní agent není povinen věřit, co mu kdo řekne (rozdíl oproti volání metody).

\subsection{KQML a~FIPA-ACL}

\begin{defbox}[KQML --- \emph{Knowledge Query and Manipulation Language}]
90.~léta, projekt DARPA \emph{Knowledge Sharing Effort} (KSE), spolu s~KIF.
\begin{itemize}
\item \textbf{KQML} = \emph{vnější} jazyk (\uv{obálka dopisu}); \emph{ilokuční} část zprávy.
\item \textbf{KIF} = \emph{vnitřní} jazyk; propoziční obsah, reprezentace znalostí.
\end{itemize}
Zpráva = \emph{performativ} + pojmenované parametry:
\begin{quote}\small\ttfamily
(ask-one\\
\hspace*{1em}:content (PRICE IBM ?price)\\
\hspace*{1em}:receiver stock-server\\
\hspace*{1em}:language LPROLOG\\
\hspace*{1em}:ontology NYSE-TICKS)
\end{quote}
\textbf{Parametry:} content, force, reply-with, in-reply-to, sender, receiver, language, ontology.

\textbf{Performativa:} achieve, advertise, ask-about, ask-one, ask-all, ask-if,
break, sorry, error, broadcast, forward, recruit-one, recruit-all, reply, subscribe, tell, \dots
\end{defbox}

\textbf{Kritika KQML:} chybí performativ \emph{commit} (závazek); nepřesná sémantika;
nesystematická a~zbytečně velká množina performativ; nestandardizovaný transport.

\begin{defbox}[FIPA-ACL --- \emph{Agent Communication Language}]
FIPA (\emph{Foundation for Intelligent Physical Agents}, dnes standardizační komise IEEE):
ACL = \textbf{zjednodušení a~zpřesnění KQML} --- menší, systematičtější množina performativ,
\emph{formálně definovaná sémantika}. Implementace: platforma JADE.

\textbf{Struktura zprávy} (13 polí --- performativ + 12 parametrů; povinné jen
\texttt{performative}):
\begin{center}\small
\begin{tabular}{ll}
\toprule
\textbf{Pole} & \textbf{Význam} \\
\midrule
\texttt{performative} & typ komunikačního aktu (ilokuční síla) --- \emph{povinné} \\
\texttt{sender}, \texttt{receiver}, \texttt{reply-to} & identifikátory agentů (AID) \\
\texttt{content} & vlastní obsah zprávy \\
\texttt{language} & jazyk obsahu (SL, KIF, \dots) \\
\texttt{encoding}, \texttt{ontology} & kódování a~ontologie pro interpretaci obsahu \\
\texttt{protocol} & interakční protokol, jehož je zpráva součástí \\
\texttt{conversation-id}, \texttt{reply-with}, \texttt{in-reply-to} & párování zpráv do konverzace \\
\texttt{reply-by} & termín odpovědi \\
\bottomrule
\end{tabular}
\end{center}
\begin{quote}\small\ttfamily
(inform\\
\hspace*{1em}:sender agent1\\
\hspace*{1em}:receiver agent2\\
\hspace*{1em}:content (price good2 150)\\
\hspace*{1em}:language sl\\
\hspace*{1em}:ontology hpl-auction)
\end{quote}
\end{defbox}

\textbf{Performativa FIPA-ACL} (22 celkem, nejdůležitější):
\begin{center}\small
\begin{tabular}{p{0.24\textwidth} p{0.70\textwidth}}
\toprule
\textbf{Skupina} & \textbf{Performativa} \\
\midrule
předávání informace & \texttt{inform}, \texttt{inform-if}, \texttt{inform-ref},
  \texttt{confirm}, \texttt{disconfirm} \\
žádost o~informaci & \texttt{query-if}, \texttt{query-ref}, \texttt{subscribe} \\
žádost o~akci & \texttt{request}, \texttt{request-when}, \texttt{request-whenever},
  \texttt{cancel} \\
vyjednávání & \texttt{cfp} (\emph{call for proposals}), \texttt{propose},
  \texttt{accept-proposal}, \texttt{reject-proposal} \\
provádění akcí & \texttt{agree}, \texttt{refuse}, \texttt{failure} \\
chyby, směrování & \texttt{not-understood}, \texttt{proxy}, \texttt{propagate} \\
\bottomrule
\end{tabular}
\end{center}

\noindent\mimoq{formální sémantika SL ve slidech není, performativa ano}

\begin{defbox}[Sémantika FIPA-ACL: FP a~RE]
Jazyk SL (\emph{Semantic Language}) --- modální logika s~operátory $B_i\varphi$ ($i$ věří),
$U_i\varphi$ (je nejistý), $I_i\varphi$ (má záměr), $\mathrm{Feasible}$, $\mathrm{Done}$.
Ke každému performativu:
\begin{itemize}
\item \textbf{FP} (\emph{feasibility precondition}) --- co musí platit, aby odesílatel mohl
  akt provést;
\item \textbf{RE} (\emph{rational effect}) --- efekt, kvůli němuž akt provádí; příjemce
  \emph{není povinen} RE naplnit (autonomie) --- RE je jen důvod odesílatele.
\end{itemize}
Pro $\langle i, \mathrm{inform}(j,\varphi)\rangle$:
\[
\text{FP: } B_i\varphi \wedge \neg B_i\big(\mathit{Bif}_j\varphi \vee \mathit{Uif}_j\varphi\big),
\qquad
\text{RE: } B_j\varphi .
\]
Čteno: \uv{věřím $\varphi$ a~nevěřím, že $j$ už o~$\varphi$ ví ($\mathit{Bif}$ --- ví, zda
platí) ani o~tom má nejisté mínění ($\mathit{Uif}$); chci, aby $j$ uvěřil $\varphi$.}
Pro $\langle i, \mathrm{request}(j,\alpha)\rangle$:
$\text{FP: } \mathit{FP}(\alpha)[i \backslash j] \wedge B_i \mathrm{Agent}(j,\alpha)
\wedge \neg B_i I_j \mathrm{Done}(\alpha)$, $\text{RE: } \mathrm{Done}(\alpha)$.

\textbf{Praktický problém:} sémantika neverifikovatelná --- zvenčí nelze ověřit, že agent
věří tomu, co tvrdí (\emph{sincerity assumption}); slabina mentalistické sémantiky ACL.
Alternativa: \emph{sociální} sémantiky nad veřejnými závazky (\emph{commitments}).
\end{defbox}

\subsection{Interakční protokoly}

Zpráva sama nestačí --- agenti vedou \textbf{konverzace} s~předepsanou strukturou.
Interakční protokol (\emph{interaction protocol}) = vzor přípustných posloupností zpráv;
FIPA standardizuje kolem tuctu (specifikace FIPA-IP, notace AUML).

\begin{center}\small
\begin{tabular}{p{0.26\textwidth} p{0.68\textwidth}}
\toprule
\textbf{Protokol} & \textbf{Průběh} \\
\midrule
\textbf{FIPA-Request}
  & $I \to P$: \texttt{request}. $P \to I$: \texttt{refuse} (konec), nebo \texttt{agree};
    po vykonání \texttt{inform-done} / \texttt{inform-result}, případně \texttt{failure}.
    Možná odpověď i~\texttt{not-understood}. \\
\textbf{FIPA-Query}
  & $I \to P$: \texttt{query-if} / \texttt{query-ref}. $P$: \texttt{refuse} nebo \texttt{agree},
    dále \texttt{inform} s~výsledkem nebo \texttt{failure}. \\
\textbf{FIPA-Contract-Net}
  & viz níže: \texttt{cfp} $\to$ $n\times$(\texttt{propose} / \texttt{refuse}) $\to$
    \texttt{accept-proposal} vítězi + \texttt{reject-proposal} ostatním $\to$
    \texttt{inform-done} / \texttt{failure}. \\
\textbf{FIPA-Iterated-Contract-Net}
  & jako CNET, iniciátor může po nabídkách vyhlásit \emph{další kolo} \texttt{cfp}
    (vyjednávání o~lepších podmínkách). \\
\textbf{FIPA-Subscribe}
  & $I$: \texttt{subscribe}; $P$: \texttt{inform} při každé změně sledované hodnoty,
    dokud nepřijde \texttt{cancel}. \\
\textbf{FIPA-Brokering / Recruiting}
  & $I$ požádá zprostředkovatele (\emph{broker}); ten najde poskytovatele, deleguje
    požadavek a~vrátí výsledky (příp.\ nechá odpovědět přímo). \\
\textbf{Aukční protokoly}
  & FIPA-English-Auction, FIPA-Dutch-Auction (viz odst.~\ref{ssec:aukce}). \\
\bottomrule
\end{tabular}
\end{center}

\textbf{Poznámka k~návrhu.} Protokol určuje legální zprávy v~daném stavu konverzace ---
\emph{předvídatelnost} autonomního chování; agent implementovatelný jako konečný automat nad
stavy konverzace (JADE: třídy \texttt{AchieveREInitiator}, \texttt{ContractNetResponder} apod.).

\subsection{Shrnutí ke~zkoušce}

\begin{itemize}
\item Ontologie = explicitní formalizovaný popis části reality (glosář + tezaurus);
  Gruberova definice. Formalismus: třídy, instance, vlastnosti, relace (tranzitivní
  \emph{is-a}), fakta. Ontologie + fakta = znalostní báze.
\item Škála formalizace od slovníku po libovolná logická omezení; aplikační / doménové /
  horní ontologie.
\item RDF = trojice \emph{subjekt--predikát--objekt}, jednoduché a~rychlé; RDFS přidává třídy.
\item OWL = deskripční logika (rozhodnutelný fragment FOL), \emph{open world assumption},
  T-Box vs.\ A-Box; profily Lite/DL/Full ($\mathcal{SHIF}$, $\mathcal{SHOIN}$),
  OWL~2 = $\mathcal{SROIQ}$ + profily EL/QL/RL.
\item KIF = FOL v~LISPovské syntaxi, obsah zprávy; KQML = obálka.
\item Epistemická logika: $K_i\varphi$ + možné světy; znalost = S5 (axiom~T), přesvědčení
  = KD45. Společná znalost $C_G\varphi$ nevznikne nespolehlivou komunikací (koordinovaný útok).
\item Komunikace agentů = \emph{akce}, ne volání metody --- příjemce není povinen vyhovět.
\item Austin: řečové akty; lokuce / ilokuce / perlokuce. Searle: podmínky zdařilosti,
  pět kategorií (asertiva, direktiva, komisiva, expresiva, deklarace); performativní sloveso
  + propoziční obsah.
\item Cohen \& Perrault: sémantika řečových aktů ve stylu STRIPS (pre / efekt) s~modálními
  operátory; umět Request(S,H,A) a~Inform(S,H,$\varphi$) --- efekt Inform =
  \emph{$H$ věří, že $S$ věří $\varphi$}.
\item KQML (obálka) + KIF (obsah); parametry a~performativa; kritika (chybí commit, nejasná
  sémantika).
\item FIPA-ACL: struktura zprávy (performative, sender/receiver, content, language, ontology,
  protocol, conversation-id, reply-by), 22 performativ, sémantika FP a~RE v~jazyce SL;
  neverifikovatelnost.
\item Protokoly: FIPA-Request, FIPA-Query, Contract Net (+ iterovaný), Subscribe, Brokering,
  aukční protokoly.
\end{itemize}

%=====================================================================
\section{Distribuované řešení problémů, kooperace, teorie her a~aukce}
\label{sec:kooperace}
%=====================================================================

Nejobsáhlejší věta okruhu, dvě situace $\Rightarrow$ dvě poloviny kapitoly:
\begin{itemize}
\item odst.\ 5.1--5.6: agenti \emph{sdílejí cíl}, jen dělba práce --- \emph{distribuované
  řešení problémů}, \emph{kooperace};
\item odst.\ 5.7--5.13: agenti s~\emph{vlastními zájmy} --- teorie her (Nashova ekvilibria,
  Paretova efektivita), mechanismy odolné vůči manipulaci (alokace zdrojů, aukce).
\end{itemize}

\subsection{Koherence a~koordinace}

Agenti: různé cíle, autonomní, pracují v~čase, rozhodují za běhu $\Rightarrow$ nutná
\emph{dynamická koordinace}. Sdílejí \textbf{úlohy} (\emph{task sharing})
a~\textbf{informace} (\emph{result sharing}).

\begin{defbox}[Koherence a~koordinace]
\begin{itemize}
\item \textbf{Koherence} (\emph{coherence}) --- kvalita systému \emph{jako celku}: kvalita
  řešení, využití zdrojů, degradace při poruchách.
\item \textbf{Koordinace} (\emph{coordination}) --- \emph{minimalizace režie} synchronizace:
  bez duplicitní práce, konfliktů o~zdroje, zbytečné komunikace.
\item Koherence možná i~při špatné koordinaci (práce třikrát, výsledek správný) --- plýtvání.
\end{itemize}
\end{defbox}

\begin{defbox}[CDPS --- \emph{Cooperative Distributed Problem Solving}]
Lesser a~kol., 80.~léta. Kooperace na problému přesahujícím individuální schopnosti
(informace, zdroje, výpočetní kapacita).
\begin{itemize}
\item Předpoklad: agenti \emph{implicitně sdílejí společný cíl}, bez konfliktů ---
  \textbf{benevolence}.
\item Měřítko úspěchu: \emph{výkon celku}; agent pomáhá celku i~ve svůj neprospěch.
\item Benevolence \emph{enormně zjednodušuje} návrh.
\end{itemize}
\end{defbox}

\begin{center}\small
\begin{tabular}{p{0.14\textwidth} p{0.38\textwidth} p{0.40\textwidth}}
\toprule
& \textbf{Charakteristika} & \textbf{Rozdíl} \\
\midrule
\textbf{PPS} \newline (\emph{parallel problem solving})
  & Bond, Gasser, 80.~léta; paralelní řešení, homogenní jednoduché procesory
  & agenti nejsou autonomní --- v~podstatě paralelní algoritmus \\
\textbf{CDPS}
  & heterogenní agenti se sofistikovanými schopnostmi, sdílený cíl, benevolence
  & konflikty se neřeší, protože nevznikají \\
\textbf{MAS}
  & agenti s~\emph{vlastními}, \emph{nesdílenými} cíli
  & nutno řešit: proč a~jak kooperovat, identifikace a~řešení konfliktů, vyjednávání,
    smlouvání \\
\bottomrule
\end{tabular}
\end{center}

\subsection{Sdílení úloh a~sdílení výsledků}

\textbf{Obecný postup CDPS} --- tři fáze:
\begin{enumerate}
\item \textbf{Dekompozice problému} --- hierarchická, rekurzivní; \emph{jak} rozložit,
  \emph{kdo} rozkládá. Krajní varianta model ACTORS: nový agent na každý podproblém, až
  po jednotlivé instrukce.
\item \textbf{Řešení podproblémů} --- sdílení informací, případně synchronizace akcí.
\item \textbf{Syntéza řešení} --- hierarchické skládání dílčích výsledků.
\end{enumerate}

\textbf{Sdílení úloh}: nutná \emph{dohoda} (kdo co udělá), typicky Contract Net.
\textbf{Sdílení výsledků}: \emph{proaktivní} (pošlu, co se podle mě bude hodit), nebo
\emph{reaktivní} (na vyžádání).

\begin{defbox}[Přínosy sdílení výsledků]
\begin{itemize}
\item \textbf{Jistota} (\emph{confidence}) --- porovnání nezávislých řešení; shoda zvyšuje
  důvěru.
\item \textbf{Úplnost} (\emph{completeness}) --- \emph{lokální} pohledy $\Rightarrow$
  globálnější představa o~problému.
\item \textbf{Přesnost} (\emph{precision}) --- zpřesnění celkového řešení.
\item \textbf{Včasnost} (\emph{timeliness}) --- kratší čas díky distribuci.
\end{itemize}
\end{defbox}

\subsection{Contract Net}

\begin{defbox}[Contract Net Protocol (CNET)]
Smith a~Davis, 1977--1980. \textbf{Sdílení úloh pomocí smluv} podle vzoru trhu: agent
v~nouzi = \emph{manažer} (\emph{manager}), ostatní = \emph{dodavatelé}
(\emph{contractors}). Čtyři fáze:
\begin{enumerate}
\item \textbf{Rozpoznání} (\emph{recognition}) --- problém nad vlastní síly (chybí
  schopnost, zdroje, informace, nebo by řešení bylo horší).
\item \textbf{Ohlášení} (\emph{announcement}) --- typicky broadcast: \emph{popis úlohy}
  (i~spustitelný), \emph{omezení} (termín, cena, kvalita), \emph{meta-informace} (kdo se
  smí hlásit).
\item \textbf{Nabídky} (\emph{bidding}) --- zájemci se schopnostmi pošlou nabídku
  (\emph{tender}): cena / čas / kvalita.
\item \textbf{Přidělení a~realizace} (\emph{awarding, expediting}) --- manažer vybere
  vítěze, smlouva; vítěz vykoná, vrátí výsledek.
\end{enumerate}
\textbf{Vlastnosti:} jednoduchý, decentralizovaný, robustní vůči výpadkům; dodavatel může
být manažerem podúlohy $\Rightarrow$ \emph{hierarchické kaskády subkontraktů}; dynamické
rozdělení zátěže. Nejčastěji implementovaný koordinační mechanismus vůbec (FIPA:
FIPA-Contract-Net; JADE: hotové chování).

\textbf{Limity:} drahý broadcast při mnoha agentech; předpoklad benevolence (dodavatel
nelže o~schopnostech) --- jinak aukce odolné vůči manipulaci (odst.~\ref{ssec:aukce}).
\end{defbox}

\slajd{mas-cnet.png}{0.98}{Contract Net jako interakční protokol FIPA-ACL. Iniciátor rozešle
\emph{cfp} (\emph{call for proposals}) s~termínem; účastník odpoví \emph{not-understood},
\emph{refuse}, nebo \emph{propose}. Iniciátor pak \emph{reject-proposal}, nebo
\emph{accept-proposal}; vítěz hlásí \emph{failure}, \emph{inform-done}, nebo
\emph{inform-ref} s~výsledkem.}

\subsection{Systémy s~tabulí (blackboard)}

\begin{defbox}[Blackboard system (BBS)]
Vůbec první schéma kooperativního řešení problémů (HEARSAY-II, rozpoznávání řeči,
70.~léta). Sdílení výsledků přes společnou datovou strukturu --- \textbf{tabuli}
(\emph{blackboard}).

Metafora: experti kolem tabule čtou i~píší; úlohy se objevují dynamicky; kdo úlohu umí,
zapíše částečné řešení; opakuje se, dokud není na tabuli řešení celé úlohy.

\textbf{Role:}
\begin{itemize}
\item \textbf{Tabule} --- sdílená paměť, typicky několik \emph{úrovní abstrakce} (hierarchie
  cílů a~akcí); zásadní formalismus reprezentace informace.
\item \textbf{Experti} (\emph{knowledge sources}) --- řeší konkrétní typ úlohy; reagují na
  cíle na tabuli, cíl z~ní přebírají, po vybrání vykonají akci; seznam schopností
  a~efektivita.
\item \textbf{Arbitr} (\emph{arbiter}, \emph{control component}) --- vybírá, kdo smí
  k~tabuli; čistě reaktivní, nebo uvažuje plány a~maximalizuje očekávanou utilitu;
  odpovídá za řešení na vyšší úrovni.
\end{itemize}
\textbf{Nutnost vzájemného vyloučení} nad tabulí --- úzké hrdlo systému.
\end{defbox}

\textbf{Klady a~zápory BBS:}
\begin{itemize}
\item[$+$] jednoduchá koordinace, kooperace, sdílení úloh i~výsledků;
\item[$+$] experti o~sobě \emph{nemusí vědět} (\emph{loose coupling}); nový expert bez
  změny ostatních;
\item[$+$] zprávy na tabuli lze přepisovat --- delegace úloh, podúlohy, změna expertů;
\item[$-$] všichni potřebují přístup k~tabuli a~stejnou architekturu; kolem tabule
  \uv{narváno} (částečně řeší distribuované hašovací tabulky).
\end{itemize}

\textbf{Příklady.}
\begin{itemize}
\item \emph{BBWar}: tabule mapuje požadované schopnosti na úlohy. \emph{Velitel}:
  \uv{dobyj město} (\texttt{ATTACK-\allowbreak CITY}) $\to$ několik \uv{zaútoč na
  pozici} (\texttt{ATTACK-\allowbreak LOCATION}). Vojáci: výběr mise podle schopností.
\item \emph{FELINE} (Wooldridge, Jennings, 1990): distribuovaný expertní systém před ACL;
  agent zveřejňuje \emph{Skills} (co umí dokázat) a~\emph{Interests} (co ho zajímá);
  komunikace trojicemi (odesílatel, příjemce, hypotéza + řečový akt).
\end{itemize}

\subsection{Distribuované omezující podmínky (DisCSP a~DCOP)}

\noindent\mimoq{není ve slidech; okruh ale mluví o~\uv{distribuovaném řešení problémů},
kam formální větev patří}

\begin{defbox}[DisCSP a~DCOP]
\begin{itemize}
\item \textbf{Distribuovaný CSP}: proměnné $x_1,\dots,x_n$ s~doménami $D_i$ rozděleny mezi
  agenty (typicky jedna na agenta); omezení mezi proměnnými \emph{různých} agentů. Cíl:
  přiřazení splňující všechna omezení. Nikdo nevidí celý problém; komunikace jen se
  sousedy v~grafu omezení.
\item \textbf{DCOP} --- optimalizační varianta: \emph{měkká} omezení jako funkce ceny
  $f_{ij}(x_i,x_j)$, maximalizuje se $\sum_{ij} f_{ij}$ = sociální blahobyt.
\item \textbf{Algoritmy:} \emph{ABT} (asynchronní backtracking, Yokoo) --- uspořádaní
  agenti, zprávy \texttt{ok?} (přiřazení níž) a~\texttt{nogood} (nahoru při konfliktu);
  dále \emph{ADOPT}, \emph{DPOP}, \emph{Max-Sum}.
\item Aplikace: rozvrhování schůzek, přidělování úkolů senzorům, alokace frekvencí.
\end{itemize}
\end{defbox}

\subsection{Interakce agentů v~prostředí}

Pomoc i~překážení i~\emph{bez komunikace}, jen přes prostředí:
\begin{itemize}
\item pozitivní interakce --- roboti přenesou cihlu, jen když tlačí ze dvou stran (nutná
  synchronizace);
\item negativní interakce --- roboti natlačení ke dveřím je nemohou otevřít (nutná
  koordinace).
\end{itemize}
\textbf{Bez znalosti typů interakcí} v~daném MAS nelze navrhnout efektivní řídicí
mechanismy. Nejjednodušší model: dva racionální (sobečtí) agenti v~prostředí připomínajícím
hru.

\subsection{Sobecký agent}

Proč být upřímný o~schopnostech? Proč dokončit přidělenou úlohu? Homogenní systém (celý
náš návrh) $\Rightarrow$ benevolence dobrá strategie. Většinou ale různé zájmy ---
konflikt společného cíle a~cílů jednotlivců (řízení letového provozu: všichni chtějí
bezpečnost, každá aerolinka chce přistát první); někdy společný zájem ani není zřejmý
$\Rightarrow$ \textbf{sobecké} (\emph{self-interested}) agenty.

\begin{defbox}[Sobecký agent]
Množina výsledků $O = \{o_1, o_2, \dots\}$ \emph{společná}, preference nikoli: agent $i$ má
vlastní \textbf{utilitní funkci} $u_i : O \to \mathbb{R}$, která výsledky uspořádá.

\textbf{Strategické myšlení:} maximalizuj svou utilitu za předpokladu racionality všech
ostatních. \emph{Nemaximalizuje} společnou utilitu, ale je \emph{robustní}.
\end{defbox}
\subsection{Hra v~normální formě}

\begin{defbox}[Hra dvou hráčů v~normální formě]
Agenti $i$, $j$ s~množinami akcí $A_i$, $A_j$; výsledek určen oběma: $e : A_i \times A_j \to O$.
\textbf{Strategie} $s_i$:
\begin{itemize}
\item \textbf{ryzí} (\emph{pure}) --- deterministická volba jedné akce;
\item \textbf{smíšená} (\emph{mixed}) --- rozdělení pravděpodobnosti nad ryzími (náhodný výběr).
\end{itemize}
\textbf{Výplatní matice} (\emph{payoff matrix}): řádky = strategie $i$, sloupce = strategie $j$,
buňka = dvojice $u_i / u_j$.
Obecně: trojice $\langle N, (A_i)_{i \in N}, (u_i)_{i \in N}\rangle$, $N$ = množina hráčů.
\end{defbox}

\begin{defbox}[Dominance]
$s_i$ \textbf{silně dominuje} $s_i'$: \emph{striktně} lepší výsledek proti \emph{všem} strategiím
protihráče. \textbf{Slabě dominuje}: alespoň stejně dobrý proti všem, striktně lepší alespoň proti
jedné (definice přednášky). \textbf{Dominantní} strategie: dominuje všem ostatním strategiím $i$.
\begin{itemize}
\item Podstatný rozdíl: pravdomluvnost ve Vickreyově aukci je jen \emph{slabě} dominantní
  (odst.~\ref{ssec:aukce}).
\item Racionální agent hraje dominantní strategii, existuje-li.
\item \emph{Iterované odstraňování dominovaných strategií} (IESDS) --- základní řešicí technika:
  dominovanou strategii nikdo nezahraje $\Rightarrow$ vyškrtnout z~matice $\Rightarrow$ další
  strategie se mohou stát dominovanými. \emph{Silně} dominované: pořadí nerozhoduje, výsledek
  jednoznačný. \emph{Slabě} dominované: na pořadí \textbf{záleží}, lze přijít o~ekvilibria.
\end{itemize}
\end{defbox}

\begin{defbox}[Nashovo ekvilibrium]
$(s_1^*, s_2^*)$ je \textbf{Nashovo ekvilibrium} (NE), pokud
\begin{itemize}
\item hraje-li $i$ strategii $s_1^*$, je pro $j$ nejlepší $s_2^*$, a~zároveň
\item hraje-li $j$ strategii $s_2^*$, je pro $i$ nejlepší $s_1^*$
\end{itemize}
--- \emph{vzájemně nejlepší odpovědi} (\emph{mutual best responses}). Pro $n$ hráčů: profil
$s^* = (s_1^*,\dots,s_n^*)$ je NE, pokud pro každého $i$ a~každou jeho $s_i$
\[
u_i(s_i^*, s_{-i}^*) \ge u_i(s_i, s_{-i}^*),
\]
$s_{-i}$ = strategie ostatních. \textbf{Žádný hráč nemá jednostranný důvod se odchýlit.}
\begin{itemize}
\item Naivní hledání: $m^n$ profilů ($n$ agentů, $m$ strategií).
\item Definice pro \emph{ryzí} strategie: NE v~ryzích nemusí existovat (kámen--nůžky--papír),
  nebo jich je více.
\end{itemize}
\end{defbox}

\begin{veta}[Nashova věta, 1950]
Každá hra s~konečným počtem hráčů a~strategií má alespoň jedno NE \emph{ve smíšených strategiích}.
\end{veta}

\textbf{Výpočetní složitost:} hledání NE je \emph{totální} vyhledávací problém (řešení vždy
existuje), \textbf{PPAD-úplný} (Daskalakis, Goldberg, Papadimitriou pro $\ge 3$ hráče; Chen
a~Deng pro dva) --- \uv{o~něco méně beznadějné než NP-úplnost}.

\begin{defbox}[Paretova efektivita a~sociální blahobyt]
\textbf{Pareto-optimální} profil (\emph{Pareto efficient}): neexistuje jiný profil zlepšující
výsledek jednoho agenta bez zhoršení jiného. Ne-Pareto-optimální řešení lze zlepšit \uv{zadarmo}.

\textbf{Sociální blahobyt} (\emph{social welfare}): $\mathit{sw}(o) = \sum_{i \in N} u_i(o)$.
Maximalizace má smysl v~kooperativních scénářích (jeden tým, jeden vlastník, jedna úloha; čím
homogennější systém, tím lépe). Maximum sw je vždy Pareto-optimální, opačně ne.
\end{defbox}

\slajd{mas-pareto.png}{0.86}{Paretova mez. \emph{Vlevo:} červené body $A$--$H$ Pareto-optimální
--- víc jedné položky jen za cenu druhé; šedé body uvnitř (např.\ $K$) nikoli --- existuje bod
lepší v~obou souřadnicích. \emph{Vpravo:} totéž pro kriteriální funkce $f_1$, $f_2$ --- $A$
dominuje $C$, mez = množina nedominovaných řešení.}
\subsection{Klasické hry}

\begin{defbox}[Vězňovo dilema (\emph{Prisoner's dilemma})]
\begin{center}
\begin{tabular}{c|cc}
$i \backslash j$ & $C$ (spolupráce) & $D$ (zrada) \\
\hline
$C$ & $3/3$ & $0/4$ \\
$D$ & $4/0$ & $1/1$ \\
\end{tabular}
\end{center}
$C$ = \emph{cooperate} (držet se spolupachatelem), $D$ = \emph{defect} (udat ho, nižší trest).
Kanonické výplaty: $R$ (\emph{reward}, oboustranná spolupráce), $P$ (\emph{punishment},
oboustranná zrada), $T$ (\emph{temptation}, zradím spolupracujícího), $S$ (\emph{sucker's
payoff}, jsem zrazen). Podmínky:
\[
T > R > P > S \qquad\text{a často navíc}\qquad 2R > T + S .
\]
\begin{itemize}
\item $R > P$: oboustranná spolupráce lepší než oboustranná zrada.
\item $T > R$, $P > S$: \emph{zrada je dominantní pro oba}.
\item $(D,D)$: \textbf{jediné NE} a~zároveň \textbf{jediné ne-Pareto-optimální} řešení.
\item $(C,C)$: maximum sociálního blahobytu ($3+3=6$).
\end{itemize}
\textbf{Individuální racionalita $\Rightarrow$ kolektivně nejhorší výsledek.}
\end{defbox}

\textbf{Důsledky vězňova dilematu:}
\begin{itemize}
\item \textbf{Tragédie obecní pastviny} (\emph{tragedy of the commons}): sdílený zdroj
  (pastvina, rybolov, atmosféra) vyčerpán --- každý maximalizuje svůj užitek.
\item Co znamená být racionální; jsou lidé racionální? (Experimentálně spolupracují častěji,
  než teorie předpovídá.)
\item \textbf{Stín budoucnosti} (\emph{shadow of the future}): v~\emph{iterovaném} PD se
  situace mění.
\end{itemize}

\begin{defbox}[Iterované vězňovo dilema a~Axelrodovy turnaje]
\begin{itemize}
\item Konečně mnoho kol a~oba to vědí: zpětná indukce $\Rightarrow$ $D$ v~každém kole
  (v~posledním není důvod spolupracovat, tedy ani v~předposledním, \dots).
\item Nekonečně / neznámý konec (pravděpodobnost pokračování): spolupráce racionální ---
  \emph{folk theorem}: každý výplatní profil lepší než minimax lze podepřít jako ekvilibrium
  opakované hry.
\end{itemize}
Axelrodovy turnaje pro iterované PD (1980): dvakrát vyhrála nejjednodušší přihlášená strategie
\textbf{Tit For Tat} (TFT, \uv{půjčka za oplátku}) --- \emph{první kolo spolupracuj, pak opakuj
poslední tah soupeře}. Vlastnosti úspěšných strategií:
\begin{itemize}
\item \textbf{slušnost} (\emph{nice}) --- nezrazuj první;
\item \textbf{odvetnost} (\emph{retaliating}) --- na zradu ihned odpověz;
\item \textbf{odpouštění} (\emph{forgiving}) --- po návratu ke spolupráci odpusť;
\item \textbf{nezávistivost} (\emph{non-envious}) --- neporážej \emph{konkrétního} soupeře;
  TFT nikdy nezískal víc bodů než protihráč, přesto vyhrál;
\item pátá: \textbf{srozumitelnost} (\emph{clarity}) --- čitelná, jinak nelze navázat
  spolupráci.
\end{itemize}
Slabina TFT: šum (chybný tah) $\Rightarrow$ dva TFT uvíznou ve vzájemné odvetě $\Rightarrow$
varianty \emph{Generous TFT} (občas odpustí), \emph{Tit for Two Tats}.
\end{defbox}
\begin{defbox}[Koordinační hra]
\begin{center}
\begin{tabular}{c|cc}
$i \backslash j$ & Balet & Zápas \\
\hline
Balet & $1/2$ & $0/0$ \\
Zápas & $0/0$ & $2/1$ \\
\end{tabular}
\end{center}
Též \emph{Battle of the Sexes} (BoS), \emph{Bach or Stravinsky}. Výplaty podporují \emph{shodu}
strategií; hráči se liší v~tom, na které shodě jim záleží víc.
\begin{itemize}
\item Ryzí NE: (Balet, Balet), (Zápas, Zápas) --- zároveň maxima sw i~Pareto-optima.
\item \textbf{Smíšené NE}: každý hraje \emph{preferovanější} možnost s~$2/3$, druhou s~$1/3$;
  očekávaná utilita jen $2/3$ --- \emph{méně} než v~kterémkoli ryzím NE.
\end{itemize}
\emph{Výpočet smíšeného NE (recept pro $2\times2$):} $j$ volí pravděpodobnost $q$ Baletu tak, aby
byl $i$ \emph{indiferentní}:
\[
u_i(\text{Balet}) = 1\cdot q + 0\cdot(1-q), \qquad u_i(\text{Zápas}) = 0\cdot q + 2\cdot(1-q).
\]
$q = 2 - 2q \Rightarrow q = 2/3$: $j$ hraje Balet (svou preferenci) s~$2/3$; symetricky $i$ hraje
Zápas s~$2/3$. Očekávaná utilita $i$: $1 \cdot \frac{2}{3} = \frac{2}{3}$.
\textbf{Klíčové pravidlo: pravděpodobnosti volím tak, aby byl \emph{protihráč} indiferentní.}
\end{defbox}

\noindent\textbf{Antikoordinační hra} (\emph{Chicken}, \emph{Hawk--Dove}): výplaty podporují
\emph{různé} strategie:
$u(\text{uhnu},\text{uhnu}) = 5/5$, $u(\text{uhnu},\text{nedám se}) = 1/6$,
$u(\text{nedám se},\text{uhnu}) = 6/1$, $u(\text{nedám se},\text{nedám se}) = 0/0$.
Ryzí NE: obě \emph{asymetrické} dvojice; maximum sw (oba uhnou, $5+5$) \emph{není} NE.
Symetrické smíšené NE: $p = 1/2$, očekávaná utilita $3$.

\subsection{Sociální volba a~volby}
\label{ssec:volby}

\noindent\mimoq{ve slidech je celá kapitola, oficiální znění okruhu ji ale nejmenuje ---
zařazeno jako doplněk k~agregaci preferencí}

Teorie her: chování agentů při daných preferencích. \emph{Teorie sociální volby} opačně:
z~individuálních preferencí \emph{odvodit} kolektivní rozhodnutí. \textbf{Funkce sociálního
blahobytu} agreguje preference do úplného uspořádání $>_{sw}$; \textbf{funkce sociální volby}
vybírá jednoho vítěze.

\textbf{Volební systémy:}
\begin{itemize}
\item \emph{většinový} --- nad 50~\%;
\item \emph{prostá většina} (\emph{plurality}) --- bod za každé první místo;
\item \emph{prostá většina s~vyřazováním} (IRV) --- opakovaně vyřaď nejslabšího, přerozděl jeho
  hlasy;
\item \emph{sekvenční většinové volby} --- pavouk párových soubojů; \emph{pořadí zápasů
  ovlivňuje výsledek};
\item \emph{Bordova metoda} --- $N-1$ bodů za první místo, \dots, $0$ za poslední; konsenzuální,
  nikoli většinové hlasování;
\item \emph{Slaterův systém} --- acyklické uspořádání minimalizující počet obrácených hran
  v~turnajovém grafu; NP-těžké.
\end{itemize}
Vítěz prosté většiny nemusí být většinový: $o_1 = 40~\%$, $o_2 = 30~\%$, $o_3 = 30~\%$
$\Rightarrow$ vyhrává $o_1$, ač ho 60~\% voličů nechtělo $\Rightarrow$ \textbf{taktické
hlasování}.

\begin{defbox}[Condorcetův paradox]
Ať zvolíme kterýkoli výsledek, většina voličů je nespokojena.
\emph{Příklad:} $V_1: A > B > C$, \quad $V_2: B > C > A$, \quad $V_3: C > A > B$.
Párové souboje: $A$ porazí $B$ (2:1), $B$ porazí $C$ (2:1), $C$ porazí $A$ (2:1) $\Rightarrow$
\textbf{kolektivní preference cyklická}, ač individuální jsou lineární; \emph{Condorcetův vítěz}
neexistuje. Stejný motiv jako neexistence ryzího NE v~kámen--nůžky--papír.
\end{defbox}

\begin{defbox}[Kritéria volebních systémů]
\begin{itemize}
\item \textbf{Paretova podmínka} (jednomyslnost): $X$ před $Y$ u~\emph{všech} voličů
  $\Rightarrow$ $X >_{sw} Y$. \emph{Splňuje} majority a~Borda, \emph{nesplňuje} sekvenční
  většinové.
\item \textbf{Podmínka Condorcetova vítěze}: kdo porazí všechny v~párovém srovnání, vyhrává.
  \emph{Splňují} sekvenční většinové, \emph{nesplňuje} plurality ani Borda.
\item \textbf{Nezávislost na irelevantních alternativách} (IIA): $X >_{sw} Y$ závisí jen na
  relativním pořadí $X$ a~$Y$ u~voličů. \emph{Nesplňuje} prakticky žádný běžný systém.
\item \textbf{Neomezený definiční obor} (univerzalita), \textbf{nediktátorství}.
\end{itemize}
\end{defbox}

\begin{veta}[Arrowova věta o~nemožnosti, 1951]
Pro tři a~více kandidátů neexistuje preferenční volební systém, který by převedl individuální
uspořádání na úplné a~tranzitivní společenské uspořádání a~zároveň splňoval neomezený definiční
obor, nediktátorství, Paretovu podmínku a~IIA.
\end{veta}

Jednodušeji: pro $>2$ kandidáty je jediná procedura s~Paretovou podmínkou a~IIA
\emph{diktatura}. \emph{Negativní} výsledek o~limitech kolektivního rozhodování: neplyne, že
jediný funkční systém je diktatura, ale že každý reálný systém některé kritérium obětuje
(nejčastěji IIA).

\textbf{Gibbardova--Satterthwaiteova věta:} každá surjektivní nediktátorská funkce sociální
volby pro $\ge 3$ kandidáty je \textbf{manipulovatelná}. Pro MAS zásadní: agent = počítač,
taktické hlasování mu nedělá morální problém $\Rightarrow$ hledají se mechanismy odolné vůči
manipulaci --- téma aukcí.
\subsection{Aukce jako mechanismus alokace zdrojů}
\label{ssec:aukce}

\begin{defbox}[Aukce]
\textbf{Tržní instituce}: zprávy obchodníků nesou cenovou informaci --- nabídka ke koupi
(\emph{bid}) nebo k~prodeji (\emph{ask}) za danou cenu; přednost mají vyšší nabídky ke koupi
a~nižší k~prodeji.
\begin{itemize}
\item V~MAS především \textbf{mechanismus rozdělení vzácných zdrojů mezi agenty}.
\item \emph{Reálně:} eBay, Sotheby's, těžební povolení, frekvenční pásma pro operátory.
  \emph{V~informatice:} procesorový čas, šířka pásma, výpočetní úlohy, reklamní pozice.
\item \textbf{Efektivita aukce}: zdroj získá agent s~nejvyšším oceněním. Prodávající
  \emph{maximalizuje} cenu, kupující \emph{minimalizuje}; každý kupující má vlastní utilitu.
\end{itemize}
\end{defbox}

\begin{defbox}[Klasifikace aukcí]
\textbf{Podle ocenění předmětu:}
\begin{itemize}
\item \emph{soukromá hodnota} (\emph{private value}) --- závisí jen na kupujícím (lístek na
  koncert);
\item \emph{společná hodnota} (\emph{common value}) --- pro všechny stejná, ale neznámá (těžební
  práva) $\Rightarrow$ \emph{prokletí vítěze} (\emph{winner's curse}): vyhrává ten, kdo se
  nejvíc spletl směrem nahoru;
\item \emph{korelovaná hodnota} --- kombinace.
\end{itemize}
\textbf{Podle protokolu:} vítěz platí první / druhou cenu; otevřené vyvolávání (\emph{open cry})
vs.\ obálková metoda (\emph{sealed bid}); jedno vs.\ více kol.
\end{defbox}
\subsection{Čtyři základní typy aukcí}

\begin{center}\small
\begin{tabular}{@{}p{0.15\textwidth}
  >{\raggedright\arraybackslash}p{0.23\textwidth}
  >{\raggedright\arraybackslash}p{0.24\textwidth}
  >{\raggedright\arraybackslash}p{0.27\textwidth}@{}}
\toprule
\textbf{Aukce} & \textbf{Protokol} & \textbf{Racionální strategie} & \textbf{Poznámky} \\
\midrule
\textbf{Anglická}
  & otevřené vyvolávání, \emph{rostoucí} cena, platí se první (nejvyšší) cena
  & \emph{dominantní}: přihazovat po $\varepsilon$ až do svého ocenění, pak odejít
  & Nejběžnější (starožitnosti, umění, internet --- asi 95~\% aukcí). Kupující typicky
    přeplatí. \\
\textbf{Holandská}
  & otevřené vyvolávání, \emph{klesající} cena, první, kdo přijme, platí aktuální cenu
  & dominantní \emph{neexistuje}; čekat těsně pod své ocenění --- hazard
  & Rychlá; zboží podléhající zkáze (květiny, ryby). Oblíbená u~prodávajících. Neodhalí tržní
    cenu ani preference ostatních. \\
\textbf{FPSB} \newline (\emph{first-price sealed bid})
  & jedno kolo, uzavřené obálky, nejvyšší nabídka vyhrává a~platí svou cenu
  & dominantní \emph{neexistuje}; nabídnout \emph{méně} než ocenění (\emph{bid shading}),
    optimum závisí na odhadu ostatních
  & Nemovitosti, státní dluhopisy. Nenutí použít utilitní funkci, neodhalí tržní cenu.
    \emph{Strategicky ekvivalentní holandské.} \\
\textbf{Vickreyova} \newline (\emph{second-price sealed bid})
  & jedno kolo, uzavřené obálky, nejvyšší nabídka vyhrává, platí \emph{druhou} nejvyšší
  & \emph{dominantní}: nabídnout \textbf{pravdivě} své ocenění
  & Teoreticky nejoblíbenější; Google AdWords, filatelistické aukce.
    \emph{Ekvivalentní anglické} (soukromé hodnoty). \\
\bottomrule
\end{tabular}
\end{center}

\begin{defbox}[Proč je pravdomluvnost dominantní strategií ve Vickreyově aukci]
Ocenění $v$, nabídka $b$, nejvyšší nabídka ostatních $p$ (neznámá). Vyhraji $\iff b > p$, platím
$p$, zisk $v - p$.
\begin{itemize}
\item \textbf{$b > v$ (nadhodnocení):} oproti $b=v$ změna jen pro $v < p < b$ --- vyhraji,
  platím $p > v$, zisk $v - p < 0$. \emph{Škodí mi.}
\item \textbf{$b < v$ (podhodnocení):} změna jen pro $b < p < v$ --- prohraji, ač bych vyhrál
  se ziskem $v - p > 0$. \emph{Přijdu o~zisk.}
\item Nabídkou $b \ne v$ \emph{nelze} získat víc než s~$b = v$.
\end{itemize}
$b = v$ je (slabě) dominantní bez ohledu na ostatní. Aukce \textbf{odolná vůči manipulaci}
(\emph{strategy-proof}, \emph{incentive compatible}) a~\emph{efektivní} (vyhrává nejvyšší
ocenění).
\end{defbox}

\medskip
\noindent\textbf{Příklad --- srovnání výnosů.} Ocenění $A = 80$~Kč, $B = 60$~Kč, $C = 30$~Kč;
ideální scénář (bez dodatečné informace, bez podvádění):
\begin{center}\small
\begin{tabular}{lll}
\toprule
\textbf{Aukce} & \textbf{Vítěz} & \textbf{Cena} \\
\midrule
Anglická & $A$ & $\approx 61$ (přihazuje se, dokud neodstoupí $B$ na 60) \\
Holandská & $A$ & $80$ nebo o~něco méně (79) --- $A$ nesmí riskovat, že přijme až pozdě \\
FPSB & $A$ & $80$ (resp.\ $60+\varepsilon$, zná-li $A$ ocenění ostatních) \\
Vickreyova & $A$ & $60$ (druhá nejvyšší nabídka) \\
\bottomrule
\end{tabular}
\end{center}
Vždy získá předmět nejvyšší ocenění (aukce \emph{efektivní}); liší se cena a~odhalená informace.

\emph{Proč u~FPSB 80, když se má nabízet pod oceněním?} Scénář předpokládá \emph{nulovou}
informaci o~ostatních; míra podhodnocení závisí na přesvědčení o~ostatních $\Rightarrow$ FPSB
(a~holandská) \textbf{nemá dominantní strategii}. Zná-li $A$ ocenění ostatních, nabídne
$60+\varepsilon$ a~výnos klesne na úroveň Vickreyovy.

\begin{defbox}[Věta o~ekvivalenci výnosů (\emph{revenue equivalence theorem})]
Vickrey (1961), Myerson, Riley--Samuelson. Předpoklady: rizikově neutrální kupující, nezávislé
soukromé hodnoty ze stejného spojitého rozdělení, předmět získá nejvyšší ocenění, kupující
s~nejnižším možným oceněním má nulový očekávaný zisk $\Rightarrow$ \textbf{všechny čtyři
základní aukce dávají prodávajícímu stejný očekávaný výnos.}

Rozdíly v~praxi = porušení předpokladů: averze k~riziku (lepší FPSB/holandská), společné hodnoty
a~prokletí vítěze (lepší anglická --- odhaluje informaci), koluze kupujících (anglická
zranitelnější), náklady na komunikaci (sealed bid levnější).
\end{defbox}
\subsection{Kombinatorické aukce}

\begin{defbox}[Kombinatorická aukce]
\emph{Více předmětů současně}, nabídky na \emph{podmnožiny} (\emph{bundles}). Motivace: hodnoty
nejsou aditivní --- \emph{komplementarita} (levá a~pravá bota, sousední frekvenční bloky:
$v(\{x,y\}) > v(x)+v(y)$), \emph{substituce} (dvě letenky na tentýž let:
$v(\{x,y\}) < v(x)+v(y)$).
\begin{itemize}
\item \textbf{Problém určení vítězů} (\emph{winner determination problem}, WDP): vybrat
  disjunktní nabídky s~maximálním součtem cen. Zobecnění množinového pakování ---
  \textbf{NP-těžký}, i~těžko aproximovatelný. Řešení: ILP solvery, větvení a~meze (CABOB, CPLEX).
\item \emph{Reprezentační} problém: $2^m - 1$ podmnožin $\Rightarrow$ kompaktní jazyky nabídek
  (OR, XOR, OR-of-XOR).
\end{itemize}
\end{defbox}
\subsection{Další typy aukcí a~alokace zdrojů}

\begin{itemize}
\item \emph{All-pay auction} --- platí všichni, vyhrává nejvyšší nabídka; model lobbingu
  a~patentových závodů.
\item \emph{Tichá} --- písemná varianta anglické; \emph{amsterodamská} --- anglická, která se
  při dvou zbylých zájemcích přepne na holandskou.
\item \emph{Dvojitá aukce} --- nabídky kupujících i~prodávajících se párují; model burzy.
\end{itemize}

\subsection{Návrh mechanismů, VCG a~vyjednávání}

\noindent\mimoq{ve slidech není; okruh jmenuje \uv{alokaci zdrojů}, k~níž toto patří ---
kombinatorické aukce slidy zmiňují jednou větou}

\begin{defbox}[Návrh mechanismů (\emph{mechanism design})]
Teorie her: \emph{jak se agenti zachovají při daných pravidlech}. Návrh mechanismů =
\uv{obrácená teorie her}: \emph{jaká pravidla, aby racionální sobečtí agenti sami dělali, co
chceme} --- typicky mluvili pravdu.
\begin{itemize}
\item Mechanismus = pravidlo výběru výsledku $x(\hat v)$ z~ohlášených preferencí + platební
  pravidlo $p_i(\hat v)$; užitek $u_i = v_i(x(\hat v)) - p_i(\hat v)$.
\item Žádoucí vlastnosti: \textbf{kompatibilita s~motivací} (vyplatí se ohlásit pravdu;
  nejsilnější varianta \emph{strategy-proof} --- pravdomluvnost dominantní),
  \textbf{individuální racionalita} ($u_i \ge 0$), \textbf{efektivita} (maximum sw),
  \textbf{vyrovnaný rozpočet}, \textbf{výpočetní zvládnutelnost}. Vše najednou nelze.
\item \textbf{Princip odhalení:} ke každému mechanismu existuje \emph{přímý} mechanismus
  (agenti ohlásí pravdu, výsledek stejný) $\Rightarrow$ teorie se točí kolem Vickreyho a~VCG.
  Bez plateb naráží na Gibbardovu--Satterthwaiteovu větu.
\end{itemize}
\end{defbox}

\begin{defbox}[VCG --- Vickrey--Clarke--Groves]
Zobecnění Vickreyovy aukce na kombinatorický případ: alokace $x^*$ maximalizující
$\sum_i \hat v_i(x)$; agent $i$ platí \textbf{externalitu způsobenou ostatním}:
\[
p_i \;=\; \max_{x} \sum_{j \ne i} \hat v_j(x) \;-\; \sum_{j \ne i} \hat v_j(x^*) .
\]
Pravdomluvný, efektivní, individuálně racionální; pro jeden předmět = Vickreyova aukce.
Nevýhody: WDP $n{+}1$-krát, malý (i~nulový) výnos prodávajícího, zranitelnost vůči koluzi.

\emph{Příklad:} předměty $\{x,y\}$; agent 1: $\{x\}\to5$, agent 2: $\{y\}\to4$, agent 3:
$\{x,y\}\to8$. Optimální alokace $5+4=9 > 8$. $p_1 = 8 - 4 = 4$ (bez agenta 1 by ostatní měli 8,
se zvolenou alokací 4), zisk $5-4=1$; symetricky $p_2 = 8-5 = 3$. Výnos prodávajícího $7 < 8$.
\end{defbox}

\begin{defbox}[Vyjednávání: monotónní protokol ústupků a~Zeuthenova strategie]
Aukce vyžaduje centrálního aukcionáře; alternativa = přímé \textbf{vyjednávání}.
\emph{Vyjednávací množina}: individuálně racionální a~Pareto-optimální dohody.
\begin{itemize}
\item \textbf{Monotónní protokol ústupků}: v~každém kole oba současně navrhnou dohodu. Konec,
  je-li návrh jednoho pro druhého alespoň tak dobrý jako jeho vlastní; jinak musí alespoň jeden
  \emph{ustoupit}, jinak vyjednávání ztroskotá.
\item \textbf{Zeuthenova strategie} --- \emph{kdo} ustoupí: ten s~menší ochotou riskovat
  konflikt:
  \[
  \mathrm{Risk}_i = \frac{u_i(\text{vlastní návrh}) - u_i(\text{návrh protějšku})}
  {u_i(\text{vlastní návrh})} .
  \]
  Ustupuje agent s~\emph{nižším} $\mathrm{Risk}$, právě tak málo, aby se poměry obrátily.
  Dvojice Zeuthenových strategií je NE; dohoda maximalizuje součin utilit (Nashovo
  vyjednávací řešení).
\end{itemize}
\end{defbox}

\subsection{Shrnutí ke~zkoušce}

\begin{itemize}
\item Koherence (výkon celku) vs.\ koordinace (minimalizace synchronizační režie).
\item CDPS: benevolence + sdílený cíl; odlišit od PPS (homogenní procesory, paralelní algoritmus)
  a~od MAS (vlastní cíle, konflikty, vyjednávání).
\item Tři fáze CDPS: dekompozice --- řešení podproblémů --- syntéza. Sdílení úloh (dohoda)
  vs.\ sdílení výsledků (proaktivní/reaktivní), čtyři přínosy: jistota, úplnost, přesnost,
  včasnost.
\item \textbf{CNET}: recognition --- announcement --- bidding --- awarding/expediting; manažer
  a~dodavatel, hierarchické subkontrakty; standardizován jako FIPA-Contract-Net.
\item \textbf{BBS}: tabule (hierarchie abstrakce, mutex), experti (skills), arbitr (výběr
  experta); volná vazba mezi experty; úzké hrdlo přístupu k~tabuli.
\item DisCSP / DCOP: proměnné a~omezení rozdělené mezi agenty, nikdo nevidí celý problém;
  ABT (\texttt{ok?} / \texttt{nogood}); DCOP maximalizuje sociální blahobyt.
\item Sobecký agent maximalizuje \emph{svoji} očekávanou utilitu za předpokladu, že ostatní
  dělají totéž; společná utilita se tím nemaximalizuje.
\item Hra v~normální formě, ryzí vs.\ smíšené strategie, dominance a~dominantní strategie,
  iterované odstraňování dominovaných strategií.
\item \textbf{Nashovo ekvilibrium} = vzájemně nejlepší odpovědi, nikdo nemá důvod se
  jednostranně odchýlit; naivní hledání $m^n$; \textbf{Nashova věta} (existence ve smíšených
  strategiích); hledání PPAD-úplné.
\item \textbf{Paretova optimalita} (nelze zlepšit jednomu bez přitížení jinému)
  vs.\ \textbf{sociální blahobyt} (součet utilit); maximum sw je Pareto-optimální, obráceně ne.
\item \textbf{Vězňovo dilema}: $T>R>P>S$; $D$ dominantní; $(D,D)$ jediné NE a~jediné
  ne-Pareto-optimální; $(C,C)$ maximalizuje sw. Tragedy of the commons.
  Iterované PD, stín budoucnosti, Axelrod a~TFT (slušná, odvetná, odpouštějící, nezávistivá,
  srozumitelná).
\item Umět spočítat smíšené NE pro $2\times2$: pravděpodobnosti volím tak, aby byl
  \emph{protihráč} indiferentní.
\item Sociální volba: plurality / IRV / sekvenční / Borda / Slater; Condorcetův paradox
  (cyklická kolektivní preference); kritéria Pareto, Condorcetův vítěz, IIA, univerzalita,
  nediktátorství; \textbf{Arrowova věta} (Pareto + IIA $\Rightarrow$ diktatura);
  Gibbard--Satterthwaite (každý systém lze manipulovat).
\item Aukce = mechanismus alokace vzácných zdrojů; efektivní aukce dá zdroj tomu, kdo ho oceňuje
  nejvíc. Dimenze: private/common value, open cry/sealed bid, first/second price, počet kol.
\item Čtyři základní typy a~strategie: anglická (přihazuj po $\varepsilon$ do svého ocenění)
  a~Vickrey (nabídni pravdivě) mají \emph{dominantní} strategii; holandská (čekej na ocenění
  $-\varepsilon$) a~FPSB (nabídni pod oceněním) dominantní strategii \textbf{nemají}.
  Umět \textbf{důkaz pravdomluvnosti} Vickreyho.
\item Ekvivalence: anglická $\leftrightarrow$ Vickreyova, holandská $\leftrightarrow$ FPSB.
  \textbf{Věta o~ekvivalenci výnosů} a~kdy neplatí (averze k~riziku, common values, prokletí
  vítěze, koluze). Umět příklad s~oceněními 80/60/30.
\item Návrh mechanismů (obrácená teorie her), princip odhalení; kombinatorické aukce a~WDP
  (NP-těžký); \textbf{VCG} = alokace maximalizující sw + platba rovná externalitě.
\item Vyjednávání: monotónní protokol ústupků + Zeuthenova strategie (ustupuje ten s~nižším
  rizikem).
\end{itemize}

%=====================================================================
\section{Metody pro učení agentů a~zpětnovazební učení}
\label{sec:uceni}
%=====================================================================

\noindent\mimoq{celá kapitola} Přednáška učení agentů výslovně \emph{nepokrývá}, okruh ho
jmenuje. Rozsah: přehled metod, MDP a~Q-learning podrobně, zbytek rámcově.
\subsection{Proč a~jak se agenti učí}

Agent v~MAS má být \emph{adaptivní}; učení = změna chování na základě zkušenosti.

\begin{defbox}[Metody pro učení agentů]
\begin{itemize}
\item \textbf{Podle učebního signálu}: \emph{s~učitelem} (správné odpovědi --- model
  protivníka z~pozorovaných tahů, klasifikace vjemů); \emph{bez učitele} (struktura ---
  shlukování situací); \textbf{zpětnovazební} (jen skalární odměna z~prostředí; pro autonomní
  agenty nejpřirozenější).
\item \textbf{Imitační učení} (\emph{learning from demonstration}): politika z~trajektorií
  experta --- \emph{behavioral cloning}, nebo \emph{inverzní RL} (rekonstrukce odměnové funkce
  z~chování experta).
\item \textbf{Evoluční metody}: populace kontrolérů šlechtěná podle utility --- typické ladění
  reaktivních architektur (kap.~\ref{sec:rizeni}); \emph{learning classifier systems},
  neuroevoluce.
\item \textbf{Podle toho, co se učí}: model prostředí, model protivníka (\emph{opponent
  modelling}), hodnotová funkce, přímo politika, u~PRS ohodnocení plánů.
\item \textbf{Individuální vs.\ sociální}: jen vlastní zkušenost, nebo i~napodobování
  a~sdílení zkušeností s~ostatními.
\end{itemize}
\end{defbox}

\textbf{Zpětnovazební učení} (\emph{reinforcement learning}, RL) --- nejběžnější,
teoreticky nejpropracovanější; pokus--omyl podle \emph{odměn} z~prostředí.
\begin{itemize}
\item Kap.~\ref{sec:architektura}: politika $\pi$ = mechanismus výběru akce, odměna
  = operacionalizace utility.
\item \textbf{Jeden agent} --- snazší, klasické RL (Q-learning).
\item \textbf{Multiagentní RL} (\emph{MARL}) --- učení \emph{současně}; formalizace Markovskými
  hrami (Nashova a~Paretova optimalita); state of the art hluboké MARL (MADRL).
\end{itemize}
\subsection{Markovský rozhodovací proces}

\begin{defbox}[MDP --- \emph{Markov Decision Process}]
Pětice $\langle S, A, T, R, \gamma\rangle$:
\begin{itemize}
\item $S$ --- stavy; $A$ --- akce;
\item $T : S \times A \times S \to [0,1]$ --- přechodová funkce,
  $T(s,a,s') = P(s_{t+1}=s' \mid s_t=s, a_t=a)$;
\item $R : S \times A \times S \to \mathbb{R}$ --- okamžitá odměna za přechod;
\item $0 \le \gamma < 1$ --- diskontní faktor (okamžitá vs.\ budoucí odměna).
\end{itemize}
\textbf{Markovská vlastnost}: další stav a~odměna závisí jen na aktuálním stavu a~akci, ne na
historii. (Přechodová funkce $\tau : R^A \to 2^E$ abstraktní architektury historii \emph{má};
MDP = její markovské zjednodušení + pravděpodobnosti + odměny.)

\textbf{Politika} (\emph{policy}) $\pi : S \to A$ (stochasticky $\pi(a|s)$). Cíl: optimální
$\pi^*$ maximalizující očekávaný diskontovaný součet odměn
\[
\mathbb{E}\Big[ \sum_{t=0}^{\infty} \gamma^t R(s_t, a_t, s_{t+1}) \Big].
\]
\textbf{Hodnotová funkce} $V^\pi : S \to \mathbb{R}$ --- očekávaná utilita stavu při politice
$\pi$; \textbf{akční hodnotová funkce} $Q^\pi(s,a)$ --- totéž pro dvojici stav--akce.
\end{defbox}

\begin{defbox}[Bellmanovy rovnice]
Pro danou politiku (\emph{Bellmanova rovnice pro evaluaci}):
\[
V^\pi(s) = \sum_{s'} T(s,\pi(s),s') \big[ R(s,\pi(s),s') + \gamma V^\pi(s') \big].
\]
Pro optimální politiku (\emph{Bellmanova rovnice optimality}):
\begin{align*}
V^*(s) &= \max_{a \in A} \sum_{s'} T(s,a,s') \big[ R(s,a,s') + \gamma V^*(s')\big], \\
Q^*(s,a) &= \sum_{s'} T(s,a,s')\big[R(s,a,s') + \gamma \max_{a'} Q^*(s',a')\big].
\end{align*}
$\pi^*(s) = \arg\max_a Q^*(s,a)$.
\end{defbox}

\noindent\textbf{Iterace hodnot} (\emph{value iteration}) --- při známém úplném popisu MDP:
\begin{algorithmic}[1]
\State inicializuj $V_0(s) \gets 0$ pro všechna $s$
\Repeat
  \ForAll{$s \in S$}
    \State $V_{k+1}(s) \gets \max_{a} \sum_{s'} T(s,a,s')\big[R(s,a,s') + \gamma V_k(s')\big]$
  \EndFor
\Until{$\max_s |V_{k+1}(s) - V_k(s)| < \varepsilon$}
\State \Return $\pi(s) = \arg\max_a \sum_{s'} T(s,a,s')[R(s,a,s') + \gamma V(s')]$
\end{algorithmic}
\begin{itemize}
\item Bellmanův operátor = kontrakce s~koeficientem $\gamma$ $\Rightarrow$ konvergence ke
  $V^*$ geometricky rychlá.
\item \textbf{Iterace politiky} (\emph{policy iteration}): střídavě \emph{vyhodnoť} politiku
  (soustava lineárních rovnic) a~\emph{zlepši} ji (hladově podle $V^\pi$); konečně mnoho
  politik $\Rightarrow$ konvergence v~konečně mnoha krocích.
\end{itemize}

\subsection{Q-learning a~SARSA}

Obvykle \textbf{neznáme} $T$ ani $R$ --- agent se učí ze zkušenosti v~diskrétních krocích:
\emph{model-free} RL.

\begin{defbox}[Q-learning (Watkins, 1989)]
\emph{Tabulka} $Q(s,a)$ --- odhad diskontovaného součtu budoucích odměn po akci $a$ ve stavu
$s$. Po přechodu $s \xrightarrow{a} s'$ s~odměnou $r$:
\[
Q(s,a) \;\leftarrow\; Q(s,a) + \alpha\Big[\, \underbrace{r + \gamma \max_{a'} Q(s',a')}
_{\text{cíl (TD target)}} - Q(s,a) \Big],
\qquad 0 \le \alpha \le 1 \text{ je rychlost učení,}
\]
ekvivalentně $Q(s,a) \leftarrow (1-\alpha)Q(s,a) + \alpha\big(r + \gamma\max_{a'}Q(s',a')\big)$.

\textbf{Off-policy:} cíl s~$\max_{a'}$ --- učí se hodnotu \emph{optimální} politiky nezávisle
na politice, kterou agent jedná. Konvergence $Q \to Q^*$: každá dvojice $(s,a)$ navštívena
nekonečněkrát, $\alpha$ vhodně klesá.
\end{defbox}

\begin{defbox}[Dilema průzkumu a~využití; $\varepsilon$-greedy]
\emph{Exploit} (využívat naučené) vs.\ \emph{explore} (zkoušet nové akce). Standardní řešení:
\[
\pi(s) = \begin{cases}
\text{náhodná akce z } A & \text{s pravděpodobností } \varepsilon,\\
\arg\max_a Q(s,a) & \text{s pravděpodobností } 1-\varepsilon .
\end{cases}
\]
$\varepsilon$ v~čase klesá (\emph{annealing}). Alternativy: softmax/Boltzmann
$\pi(a|s) \propto e^{Q(s,a)/\tau}$, optimistická inicializace $Q$, UCB (horní meze
spolehlivosti).
\end{defbox}

\medskip
\noindent\textbf{Konkrétní krok Q-learningu.} $\alpha = 0{,}5$, $\gamma = 0{,}9$; ve stavu
$s_1$ akce $a_1$, odměna $r = 10$, přechod do $s_2$. Hodnoty: $Q(s_1,a_1) = 4$,
$Q(s_2,a_1) = 6$, $Q(s_2,a_2) = 8$.
\begin{align*}
\max_{a'} Q(s_2,a') &= \max(6,\,8) = 8, \\
\text{TD cíl} &= r + \gamma \max_{a'} Q(s_2,a') = 10 + 0{,}9 \cdot 8 = 17{,}2, \\
\text{TD chyba} &= 17{,}2 - Q(s_1,a_1) = 17{,}2 - 4 = 13{,}2, \\
Q(s_1,a_1) &\leftarrow 4 + 0{,}5 \cdot 13{,}2 = \mathbf{10{,}6}.
\end{align*}

\begin{defbox}[SARSA]
\emph{State--Action--Reward--State--Action}: v~aktualizaci \emph{skutečně zvolená} další akce
$a'$ (touž politikou, kterou agent jedná):
\[
Q(s,a) \leftarrow Q(s,a) + \alpha\big[r + \gamma\, Q(s',a') - Q(s,a)\big].
\]
\textbf{On-policy:} učí se hodnotu prováděné politiky včetně náhodného průzkumu $\Rightarrow$
\emph{opatrnější}. \uv{Chůze po útesu} (\emph{cliff walking}): Q-learning --- optimální cesta
těsně u~propasti, při $\varepsilon$-greedy občas spadne; SARSA --- bezpečnější cesta dál od
okraje, započítá riziko vlastního průzkumu.

Zobecnění obou: $\mathrm{TD}(\lambda)$ s~\emph{eligibility traces} (kredit rozprostřen na více
předchozích kroků).
\end{defbox}

\subsection{Metody založené na politice}

\begin{defbox}[Policy gradient a~actor--critic]
Optimalizace \emph{přímo politiky}: parametrizace $\pi(a \mid s, z)$, hledáme $z^*$
gradientním výstupem po očekávaném výnosu. \textbf{REINFORCE} (Williams, 1992) --- odhad
gradientu z~Monte Carlo simulací celých epizod:
\[
z_{t+1} = z_t + \alpha\, G_t\, \nabla_z \ln \pi(A_t \mid S_t, z_t),
\]
$G_t$ = výnos od času $t$; intuice: zvyš pravděpodobnost akcí následovaných vysokým výnosem.
\begin{itemize}
\item Výhody proti hodnotovým metodám: \emph{spojité} akční prostory; přirozeně
  \emph{stochastické} politiky (nutné ve hrách se smíšenými ekvilibrii).
\item Nevýhoda: vysoký rozptyl odhadu gradientu.
\end{itemize}

\textbf{Actor--critic} --- kombinace: \emph{actor} = politika, \emph{critic} = hodnotová
funkce (snižuje rozptyl). \textbf{Funkce výhody} $A(s,a) = Q(s,a) - V(s)$: gradient škálován
relativní výhodností akce místo výnosu. Varianty: A3C, A2C, PPO, SAC, DDPG.

\textbf{Hluboké RL:} \textbf{DQN} --- tabulka $Q$ nahrazena neuronovou sítí; triky
\emph{experience replay} (rozbití korelací mezi vzorky), \emph{target network} (stabilizace
cíle).
\end{defbox}

\subsection{Multiagentní zpětnovazební učení}

\begin{defbox}[Markovská hra (\emph{Markov game}, stochastická hra)]
MDP pro $N$ agentů: $\langle N, S, (A_i)_{i\in N}, T, (R_i)_{i \in N}, \gamma\rangle$.
\textbf{Sdružený akční prostor} $\mathbf{A} = A_1 \times \cdots \times A_N$; přechod závisí na
akcích \emph{všech}; každý agent má \emph{vlastní} $R_i$ a~$\pi_i$, jeho hodnotová funkce
závisí na politikách všech.

\textbf{Vazba na teorii her:} politika agenta = \emph{nejlepší odpověď} na sdružené politiky
ostatních; sdružené politiky mohou být v~\textbf{Nashově ekvilibriu} a~\textbf{Pareto-optimální}.
Speciální případy: $N=1$ --- MDP; jeden stav --- hra v~normální formě; $N=2$ s~nulovým součtem
--- \emph{minimax} hra.

\textbf{Minimax-Q} (Littman, 1994) --- dva agenti, nulový součet; místo $\max_{a'} Q(s',a')$
hodnota smíšeného minimaxového ekvilibria matice $Q(s',\cdot,\cdot)$ (lineární program).
\textbf{Nash-Q} --- totéž pro obecný součet, konverguje jen za silných předpokladů.
\end{defbox}

\begin{defbox}[Proč je MARL těžké]
\begin{itemize}
\item Politiky aktualizovány \textbf{současně} $\Rightarrow$ prostředí z~pohledu agenta
  \textbf{nestacionární} (\uv{pohyblivý cíl}); porušen základní předpoklad konvergence RL.
\item Sdružený akční prostor roste \textbf{exponenciálně} s~počtem agentů.
\item Jen \textbf{částečná pozorovatelnost} (Dec-POMDP; optimální řešení NEXP-úplné).
\item \textbf{Přiřazení zásluh}: při společné odměně nejasné, kdo přispěl.
\end{itemize}
\textbf{Protiopatření:} \emph{centralizované učení s~decentralizovaným prováděním} (CTDE) ---
kritik při tréninku vidí stav i~akce všech, herec za běhu jen svá lokální pozorování (MADDPG,
COMA, QMIX); sdílení parametrů mezi homogenními agenty; \emph{opponent modelling}.
\end{defbox}

\textbf{Self-play.} Učení hrou proti sobě (či dřívějším verzím sebe); v~hrách s~nulovým součtem
spolehlivě silné strategie (TD-Gammon, AlphaZero). Riziko: \emph{cyklické} preference strategií
$\Rightarrow$ naivní self-play osciluje; proto ligy protivníků a~populační metody.

\subsection{Shrnutí ke~zkoušce}

\begin{itemize}
\item Metody učení agentů: s~učitelem / bez učitele / zpětnovazební; imitační učení
  a~inverzní RL; evoluční metody (LCS, neuroevoluce); co se agent učí (model prostředí,
  model protivníka, politiku); individuální vs.\ sociální učení. RL je nejběžnější.
\item \textbf{MDP} $= \langle S,A,T,R,\gamma\rangle$, Markovská vlastnost, politika $\pi$,
  hodnotové funkce $V^\pi$, $Q^\pi$; Bellmanova rovnice optimality; value iteration
  (kontrakce, konverguje) a~policy iteration.
\item \textbf{Q-learning}: $Q(s,a) \leftarrow Q(s,a) + \alpha[r + \gamma\max_{a'}Q(s',a') - Q(s,a)]$,
  off-policy, tabulkový; umět spočítat jeden krok s~čísly.
\item \textbf{SARSA}: $\dots + \alpha[r + \gamma Q(s',a') - Q(s,a)]$, on-policy, opatrnější
  (cliff walking). $\varepsilon$-greedy jako řešení dilematu průzkum/využití.
\item \textbf{Policy gradient} (REINFORCE, $\nabla \ln \pi$ krát výnos) a~\textbf{actor-critic}
  (herec = politika, kritik = hodnotová funkce, výhoda $A = Q - V$); DQN (replay buffer,
  target network).
\item \textbf{Markovská hra} = MDP pro $N$ agentů se sdruženým akčním prostorem
  a~individuálními odměnami; nejlepší odpověď, Nashovo ekvilibrium, Pareto-optimalita.
  \textbf{Minimax-Q} (nulový součet, LP v~každé aktualizaci), Nash-Q (obecný součet).
\item Problémy MARL: nestacionarita, exponenciální sdružený akční prostor, částečná
  pozorovatelnost, přiřazení zásluh; řešení CTDE. Self-play a~AlphaZero.
\end{itemize}

%=====================================================================
\section{Metodologie návrhu, jazyky a~prostředí MAS}
\label{sec:metodologie}
%=====================================================================

\subsection{Agentově orientované softwarové inženýrství}
Proč nestačí OO metodologie (UML):
\begin{itemize}
\item agenti \emph{autonomní} --- nelze přímo volat jejich metody;
\item vlastní cíle;
\item vztahy \emph{organizační} (role, protokoly, závazky), nikoli strukturální.
\end{itemize}
Odpověď: metodologie \textbf{AOSE} (\emph{agent-oriented software engineering}).

\begin{defbox}[Gaia (Wooldridge, Jennings, Kinny, 2000)]
Jedna z~nejcitovanějších AOSE metodologií. Metafora: MAS = \textbf{organizace}, agenti =
\textbf{role}. Postup: požadavky $\to$ analýza $\to$ návrh; implementace až na konkrétní
platformě.

\textbf{Analýza} --- dva abstraktní modely:
\begin{itemize}
\item \emph{Model rolí} (\emph{roles model}); role = čtyři atributy:
  \emph{permissions} --- zdroje, které smí číst/měnit;
  \emph{responsibilities} --- co má zajistit: \emph{liveness} (\uv{něco dobrého se stane},
  regulární výraz nad aktivitami a~protokoly) a~\emph{safety} (\uv{nic špatného se nestane},
  invarianty);
  \emph{activities} --- výpočty prováděné rolí samotnou;
  \emph{protocols} --- interakce s~jinými rolemi.
\item \emph{Model interakcí}: protokoly mezi rolemi --- účel, iniciátor, respondér, vstupy,
  výstupy.
\end{itemize}

\textbf{Návrh} --- tři konkrétní modely:
\begin{itemize}
\item \emph{model agentů}: role $\to$ typy agentů, počty instancí;
\item \emph{model služeb}: funkce poskytované rolí --- vstupy, výstupy, předpoklady, efekty;
\item \emph{model známostí} (\emph{acquaintance model}): orientovaný graf, kdo s~kým
  komunikuje --- kontrola úzkých hrdel a~přílišné provázanosti.
\end{itemize}

\textbf{Omezení:} uzavřený systém, kooperativní agenti, statická organizace; uvolňují
rozšíření \emph{ROADMAP} a~\emph{Gaia v.2}.
\end{defbox}

\subsection{Standard FIPA a~architektura platformy}

\begin{defbox}[Referenční architektura agentní platformy podle FIPA]
\begin{itemize}
\item \textbf{AMS} (\emph{Agent Management System}) --- povinná; \uv{matrika} platformy:
  evidence agentů, jednoznačné \emph{AID} (agent identifier), životní cyklus podle
  stavového diagramu FIPA (\emph{initiated} $\to$ \emph{active} $\leftrightarrow$
  \emph{waiting}/\emph{suspended}, u~mobilních agentů navíc \emph{transit}; operace create,
  suspend, resume, terminate), bílé stránky (\emph{white pages}).
\item \textbf{DF} (\emph{Directory Facilitator}) --- žluté stránky (\emph{yellow pages}):
  registrace nabízených \emph{služeb}, vyhledávání jejich poskytovatelů.
\item \textbf{MTS} (\emph{Message Transport Service}) --- doručování ACL zpráv v~rámci
  platformy i~mezi platformami (protokoly IIOP, HTTP; kódování bitefficient, XML, string).
\item \textbf{ACC} (\emph{Agent Communication Channel}) --- rozhraní k~MTS.
\end{itemize}
Standard dále definuje: jazyk ACL, obsahové jazyky (SL), interakční protokoly, správu
ontologií.
\end{defbox}

\subsection{Jazyky a~prostředí}

\noindent\mimoq{AGENT-0 a~Concurrent MetateM ve slidech nejsou; JADE a~Jason ano}
\begin{itemize}
\item \textbf{Agentově orientované programování} (Shoham 1993) --- historické východisko:
  stav programu = \emph{mentální kategorie} (přesvědčení, závazky, schopnosti), program =
  množina \emph{pravidel závazků}.
\item \textbf{AGENT-0} --- první jazyk AOP.
\item \textbf{Concurrent MetateM} --- jiná cesta: agent = \emph{přímo spustitelná
  specifikace} v~temporální logice.
\item Oba navazují na deduktivní agenty (odst.\ 2.1); prakticky nejúspěšnější agentní jazyk
  dnes: AgentSpeak/Jason.
\end{itemize}
\begin{center}\small
\begin{tabular}{p{0.15\textwidth} p{0.20\textwidth} p{0.57\textwidth}}
\toprule
\textbf{Systém} & \textbf{Typ} & \textbf{Charakteristika} \\
\midrule
\textbf{JADE}
  & Java, FIPA platforma
  & \emph{Java Agent DEvelopment Framework}; jedna z~nejrozšířenějších platforem, plně
    FIPA. Agent = třída dědící z~\texttt{Agent}, metoda \texttt{setup()};
    chování = objekty \texttt{Behaviour} (\texttt{OneShot}, \texttt{Cyclic},
    \texttt{Ticker}, \texttt{FSMBehaviour}, \texttt{SequentialBehaviour}), kooperativní
    (nepreemptivní) plánovač v~jednom vlákně agenta. Zprávy: \texttt{ACLMessage},
    \texttt{send()}/\texttt{blockingReceive()}, \texttt{MessageTemplate}. Hotové protokoly
    (\texttt{ContractNetInitiator}, \texttt{AchieveREResponder}), ladicí nástroje (RMA,
    Sniffer, Introspector), kontejnery, mobilita agentů. \\
\textbf{Jason}
  & interpret AgentSpeak(L) v~Javě
  & Referenční implementace \textbf{AgentSpeak} = \textbf{BDI} jazyk (logické programování
    + BDI), přímý potomek PRS. Program = počáteční \emph{beliefs}, \emph{rules},
    \emph{knihovna plánů} tvaru \texttt{spouštěč : kontext <- tělo}. Spouštěč =
    \emph{událost}: \texttt{+b} přidání přesvědčení, \texttt{+!g} přidání cíle, \texttt{-!g}
    selhání cíle. Cyklus interpretu: vnímej $\to$ aktualizuj beliefs $\to$ vyber událost
    $\to$ \emph{relevantní} plány (podle spouštěče) $\to$ \emph{aplikovatelné} (podle
    kontextu) $\to$ vytvoř záměr $\to$ vykonej krok. Jason + \texttt{CArtAgO} (prostředí)
    + \texttt{Moise} (organizace) = \emph{JaCaMo}. \\
\textbf{NetLogo}
  & prostředí pro ABM
  & Jazyk odvozený od Logo pro \emph{agent-based modeling}. Entity: \emph{turtles} (mobilní
    agenti), \emph{patches} (buňky prostředí), \emph{links}, \emph{observer}. Knihovna
    modelů (Segregation, Wolf Sheep Predation, Ants, Flocking), \emph{BehaviorSpace} pro
    parametrické studie. Emergentní chování a~výuka, ne FIPA komunikace. \\
\textbf{SPADE}
  & Python, XMPP
  & \emph{Smart Python Agent Development Environment}. Komunikace standardním
    \textbf{XMPP} místo vlastního MTS; agent = třída s~\emph{behaviours}
    (\texttt{CyclicBehaviour}, \texttt{PeriodicBehaviour}, \texttt{FSMBehaviour}); asyncio,
    webové rozhraní; populární díky napojení na Python ekosystém ML. \\
\bottomrule
\end{tabular}
\end{center}

\begin{defbox}[Jason podrobněji (látka cvičení NAIL106)]
\textbf{Jazyk.}
\begin{itemize}
\item Spouštěče \texttt{+b}, \texttt{-b}, \texttt{+!g}, \texttt{-!g}; navíc \emph{testovací
  cíl} \texttt{?g} --- dotaz do báze beliefs, unifikací naplní proměnné.
\item Tělo plánu mění beliefs přímo: \texttt{+b} přidání, \texttt{-b} odebrání, \texttt{-+b}
  aktualizace.
\item \textbf{Dvě negace}: \texttt{not b} = negace jako neúspěch (\uv{agent nemá důvod věřit
  \texttt{b}}); \texttt{\textasciitilde b} = silná negace (\uv{agent věří, že \texttt{b}
  neplatí}).
\item \emph{Anotace} beliefs, hlavně zdroj: \texttt{likes(a,b)[source(martin)]} --- zdroj =
  vjem (\texttt{percept}), agent sám (\texttt{self}), nebo jiný agent.
\end{itemize}

\textbf{Komunikace.} Interní akce \texttt{.send}(\emph{příjemce}, \emph{performativ},
\emph{obsah}); performativy \texttt{tell}, \texttt{untell}, \texttt{achieve},
\texttt{askOne}, \dots{} = řečové akty ve smyslu KQML/FIPA-ACL (kap.~\ref{sec:komunikace}).
\texttt{tell} přidá příjemci belief s~anotací \texttt{source(odesílatel)}, \texttt{achieve}
mu přidá cíl. Ve slidech pivní robot z~odd.~\ref{sec:rizeni}:
\texttt{.send(supermarket, achieve, order(beer,5))}.

\textbf{Prostředí a~projekt.} Prostředí = třída v~Javě dědící z~\texttt{Environment}:
přepsat \texttt{executeAction}(\emph{agent}, \emph{akce}), vjemy agentům přes
\texttt{addPercept}; pro mřížkové světy hotový \texttt{GridWorldModel}. Projektový soubor
\texttt{.mas2j}: seznam agentů (a~počty instancí), třída prostředí, infrastruktura.

\textbf{Přizpůsobení interpretu.}
\begin{itemize}
\item Tři volby cyklu (vyber událost / aplikovatelný plán / záměr) = tři \emph{selekční
  funkce} \texttt{selectEvent}, \texttt{selectOption}, \texttt{selectIntention}; přepis
  děděním z~\texttt{Agent} (např.\ prioritizace typu událostí).
\item Dědit lze celou architekturu \texttt{AgArch} (\texttt{perceive}, \texttt{checkMail}
  --- vlastní model světa z~vjemů a~zpráv).
\item Vlastní \texttt{BeliefBase} (např.\ automatická náhrada duplicitních beliefs).
\item Výpočetně těžké podprogramy (typicky hledání cest) = vlastní interní akce děděním
  z~\texttt{DefaultInternalAction}.
\end{itemize}
\end{defbox}

\textbf{BattleCode} (slidy seminářů kurzu vedle JADE a~AgentSpeak/Jasonu): programátorská
soutěž MIT --- bitevní strategie + softwarové inženýrství + UI; tým píše program autonomně
řídící každého robota své armády zvlášť (decentralizované řízení, omezený výpočet na tah,
omezená komunikace) = MAS v~malém jako cvičiště.

\textbf{Další prostředí:}
\begin{itemize}
\item \emph{Jadex}, \emph{JACK}, \emph{JAM} --- další implementace BDI/PRS (Jadex = BDI
  vrstva nad JADE).
\item \emph{MASON}, \emph{Repast}, \emph{GAMA} --- simulační platformy pro agentní modely.
\item Multiagentní zpětnovazební učení (kap.~\ref{sec:uceni}):
  \textbf{Gymnasium} (nástupce OpenAI Gym) --- standardní API jednoagentního RL:
  \texttt{reset()}, \texttt{step(akce)} $\to$ pozorování, odměna, příznak konce;
  \textbf{PettingZoo} --- multiagentní obdoba, dvě API: \emph{AEC} (agenti táhnou střídavě)
  a~\emph{Parallel} (všichni současně), hotová prostředí (MPE, multiplayer Atari, deskové
  hry);
  \textbf{OpenSpiel} (DeepMind) --- kolekce her v~extenzivní i~normální formě (nulový
  i~obecný součet, imperfektní informace) a~referenčních algoritmů (CFR, fictitious play,
  MCTS) pro učení ve hrách.
\end{itemize}

\subsection{Shrnutí ke~zkoušce}

\begin{itemize}
\item Proč nestačí OO metodologie: autonomie, vlastní cíle, organizační (nikoli strukturální)
  vztahy.
\item \textbf{Gaia}: organizace + role; analýza = model rolí (permissions, responsibilities
  s~liveness a~safety, activities, protocols) + model interakcí; návrh = model agentů,
  služeb a~známostí. Předpoklad kooperativních agentů a~statické organizace.
\item \textbf{FIPA platforma}: AMS (bílé stránky, AID, životní cyklus: initiated, active,
  waiting, suspended, transit), DF (žluté stránky, služby), MTS/ACC (doručování ACL zpráv).
\item \textbf{JADE} = Java, FIPA, \texttt{Agent} + \texttt{Behaviour} + \texttt{ACLMessage},
  hotové protokoly a~ladicí nástroje. \textbf{Jason/AgentSpeak} = BDI jazyk,
  plán \texttt{spouštěč : kontext <- tělo}, potomek PRS. \textbf{NetLogo} = ABM simulace
  (turtles/patches/links). \textbf{SPADE} = Python + XMPP.
\item \textbf{Jason prakticky}: testovací cíl \texttt{?g}, aktualizace \texttt{-+b};
  \texttt{not} (negace jako neúspěch) vs.\ \texttt{\textasciitilde} (silná negace); anotace
  \texttt{[source(...)]}; \texttt{.send} s~performativy tell/achieve = řečové akty; prostředí
  v~Javě (\texttt{executeAction}, \texttt{addPercept}), projekt \texttt{.mas2j}; přizpůsobení
  přes selekční funkce cyklu, \texttt{AgArch} a~vlastní interní akce.
\end{itemize}

\section*{Závěrem: jak na zkoušku}
\addcontentsline{toc}{section}{Závěrem: jak na zkoušku}

Zkoušející ověřuje \emph{souvislosti}. Čtyři osy otázek:
\begin{enumerate}
\item \textbf{Reaktivita vs.\ proaktivita} --- definice inteligentního agenta $\to$ commitment
  strategie (bold vs.\ cautious), parametry $\gamma/\phi$ v~síti Maes $\to$ vrstvy hybridních
  architektur.
\item \textbf{Individuální vs.\ kolektivní racionalita} --- sobecký agent $\to$ $(D,D)$ ve
  vězňově dilematu, tragédie obecní pastviny, taktické hlasování $\to$ potřeba mechanism
  designu (Vickrey, VCG): návrh \emph{pravidel}, za nichž se sobeckému agentovi vyplatí
  pravdomluvnost.
\item \textbf{Symbolická vs.\ subsymbolická reprezentace} --- FOL, STRIPS, ontologie, ACL
  vs.\ subsumpce, aktivační sítě, neuronové sítě v~RL; praxe hybridní.
\item \textbf{Teorie vs.\ praxe} --- elegantní, ale nekonstruktivní definice (optimální agent,
  Nashova věta, Arrowova věta) vs.\ použitelné algoritmy (A*, CBS, Q-learning, CNET, JADE).
  Složitost: PPAD-úplnost hledání NE, NP-těžkost WDP i~optimálního MAPF.
\end{enumerate}

Časté drobnosti: \emph{desire} vs.\ \emph{intention}; proč efektem \texttt{inform}
\emph{není}, že příjemce uvěří; proč je $(D,D)$ jediné \emph{ne}-Pareto-optimální řešení
vězňova dilematu; proč je pravdomluvnost dominantní ve Vickreyově aukci; off-policy
Q-learning vs.\ on-policy SARSA; co dělá CBS na vysoké a~co na nízké úrovni.

\end{document}
